\documentclass[11pt,twoside]{amsart}
\usepackage[dvipsnames]{xcolor}
\usepackage{hyperref}
\hypersetup{colorlinks=true, linkcolor=Blue, citecolor=Blue}
\usepackage{amsthm, amsmath, amscd, amssymb,centernot,txfonts}
\usepackage{tikz-cd}
\usepackage{lipsum}
\usepackage{cancel} 
\usepackage{multicol}
\usepackage{amsfonts}
\usepackage{amsmath}
\usepackage{mathrsfs}
\usepackage[all]{xy}
\usepackage[left=2.5cm,top=2.5cm,bottom=3cm,right=2.5cm]{geometry}
\usepackage{thmtools}
\usepackage{dirtytalk}
\usepackage{mathtools}
\usetikzlibrary{shapes,arrows}
\usepackage{etoolbox}
\usepackage{dynkin-diagrams}
\usepackage{enumitem}
\usepackage{chngpage}
\usepackage{adjustbox}
\usepackage[normalem]{ulem}
\usepackage[utf8]{inputenc}
\usepackage{fourier} 
\usepackage{array}
\usepackage{makecell}
\usepackage{comment}
%\usepackage[cmtip,all]{xy}
\usepackage{cancel}
\usepackage{cleveref}
\newcommand{\longsquiggly}{\xymatrix{{}\ar@{~>}[r]&{}}}




\setlength{\headheight}{15.2pt}

\renewcommand\theadalign{bc}
\renewcommand\theadfont{\bfseries}
\renewcommand\theadgape{\Gape[4pt]}
\renewcommand\cellgape{\Gape[4pt]}



\setcounter{tocdepth}{1}

\setlength\parskip{0in}
\setlength\parindent{0.2in}


\theoremstyle{plain}
\numberwithin{equation}{section}
\newtheorem{theorem}{Theorem}[section]
\newtheorem{proposition}[theorem]{Proposition}
\newtheorem{lemma}[theorem]{Lemma}
\newtheorem{corollary}[theorem]{Corollary}
\newtheorem{conjecture}[theorem]{Conjecture}
\newtheorem{claim}[theorem]{Claim}
%\newtheorem{set-up}[theorem]{Set-up}

\DeclareMathOperator{\Spec}{Spec}
\DeclareMathOperator{\Proj}{Proj}
%\DeclareMathOperator{\Hom}{Hom}
\DeclareMathOperator{\Der}{Der}
\DeclareMathOperator{\coker}{coker}
%\DeclareMathOperator{\Ext}{\mathcal{E}xt}
\DeclareMathOperator{\im}{im}
\DeclareMathOperator{\rk}{rk}
 
\newcommand{\OO}{\mathcal{O}}
\newcommand{\PP}{\mathbb{P}}
\newcommand{\AAA}{\mathbb{A}}
\newcommand{\CC}{\mathbb{C}}
\newcommand{\ZZ}{\mathbb{Z}}
\newcommand{\Tt}{\mathcal{T}}


\newtheorem{observation}[theorem]{Observation}


\theoremstyle{definition}
\newtheorem{remark}[theorem]{Remark}
\newtheorem{notation}[theorem]{Notation}
\newtheorem{example}[theorem]{Example}
\newtheorem{question}[theorem]{Question}
\newtheorem{set-up}[theorem]{Set-up}
\newtheorem{definition}[theorem]{Definition}
\newtheorem{construction}[definition]{Construction}
\newcommand{\te}{\otimes}

\newcommand{\myeq}{\overset{\mathrm{duality}}{=\joinrel=}}
\newcommand*{\QEDA}{\hfill\ensuremath{\blacksquare}}
\newcommand*{\QEDB}{\hfill\ensuremath{\square}}

\newcommand{\cO}{\mathscr{O}}
\newcommand{\cE}{\mathscr{E}}
\newcommand{\cB}{\mathscr{B}}
\newcommand{\cI}{\mathscr{I}}
\newcommand{\Yt}{\widetilde{Y}}
\newcommand{\Ext}{\operatorname{Ext}}
\newcommand{\Aut}{\operatorname{Aut}}
\newcommand{\Hom}{\operatorname{Hom}}

%\DeclareMathOperator{\Ext}{Ext}
%\DeclareMathOperator{\Hom}{Hom}
%\DeclareMathOperator{\coker}{coker}
\DeclareMathOperator{\Tot}{Tot}
\DeclareMathOperator{\Div}{Div}
\DeclareMathOperator{\Pic}{Pic}
\DeclareMathOperator{\Sing}{Sing}
%\DeclareMathOperator{\Spec}{Spec}
\DeclareMathOperator{\Sym}{Sym}
\DeclareMathOperator{\height}{ht}
 
%\newcommand{\OO}{\mathcal{O}}
\newcommand{\LL}{\mathbf{L}}
\newcommand{\cT}{\mathcal{T}}
\newcommand{\cH}{\mathcal{H}}
%\newcommand{\cI}{\mathcal{I}}
\newcommand{\cJ}{\mathcal{J}}
\newcommand{\cN}{\mathcal{N}}
%\newcommand{\PP}{\mathbb{P}}
%\newcommand{\CC}{\mathbb{C}}
%\newcommand{\ZZ}{\mathbb{Z}}
\newcommand{\HH}{\mathrm{H}}


\tikzstyle{decision} = [diamond, draw, , 
    text width=4.5em, text badly centered, node distance=3cm, inner sep=0pt]
\tikzstyle{block} = [rectangle, draw, , 
    text width=10em, text centered, rounded corners, minimum height=2em]
\tikzstyle{block1} = [rectangle, draw, , 
    text width=5em, text centered, rounded corners, minimum height=2em]    
\tikzstyle{line} = [draw, -latex']
\tikzstyle{cloud} = [draw, ellipse,, node distance=3cm,
    minimum height=2em]

\begin{document}

\title[Deformation theory of morphisms finite onto a singular base ]{Deformation theory of morphisms finite onto a singular base}

\author[J. Mukherjee]{Jayan Mukherjee}
\address{Department of Mathematics, Oklahoma State University, Stillwater, USA}
\email{jayan.mukherjee@okstate.edu}

\maketitle

\section{First--order, locally trivial, infinitesimal deformations of a morphism}



\begin{remark}(see \cite[2.4]{Gon06})\label{D(X F)}
The proof given in , in the context of complex manifolds, is still valid
in the case of a homomorphism of quasi--coherent sheaves
\[
\mathcal A \xrightarrow{\,F\,} \mathcal B
\]
with kernel and cokernel, respectively, $\mathcal K$ and $\mathcal N$, on a noetherian separated scheme $X$.
Let $\mathcal V$ be an open cover of $X$. We set
\begin{equation}\tag{2.4.1}
D(X,F;\mathcal V)
=
\frac{\Bigl\{(g,\rho)\in C^{0}(\mathcal V,\mathcal B)\times Z^{1}(\mathcal V,\mathcal A)\ \Big|\ 
\delta g = F\rho\Bigr\}}
{\Bigl\{(Fh,\delta h)\ \Big|\ h\in C^{0}(\mathcal V,\mathcal A)\Bigr\}}
\end{equation}
and
\begin{equation}\tag{2.4.2}
D(X,F)=\varinjlim_{\mathcal V}\, D(X,F;\mathcal V),
\end{equation}
where the direct limit is taken under refinement of coverings. Then for every open affine cover
$\mathcal V$ the natural map $D(X,F;\mathcal V)\to D(X,F)$ is an isomorphism.
Two exact sequences, like in Lemma~2.3, are then obtained:
\[
H^{0}(\mathcal A)\longrightarrow H^{0}(\mathcal B)\longrightarrow D(X,F)
\longrightarrow H^{1}(\mathcal A)\longrightarrow H^{1}(\mathcal B),
\]
\[
0\longrightarrow H^{1}(\mathcal K)\longrightarrow D(X,F)\longrightarrow H^{0}(\mathcal N)
\longrightarrow H^{2}(\mathcal K).
\]

Consequently, if $F$ is injective, $D(X,F) \cong H^0(\mathcal{N})$. 

\end{remark}

Let $X \xrightarrow{\varphi} Z$ be a morphism, where $X$ is a reduced connected scheme and $Z$
is a scheme. Let $\Delta$ denote $\operatorname{Spec} k[\epsilon]/(\epsilon^{2})$. We are interested in the
first--order, locally trivial, infinitesimal deformations of the pair $(X,\varphi)$, i.e.\ $\Delta$--morphisms
\[
\widetilde{X} \xrightarrow{\widetilde{\varphi}} Z \times \Delta
\]
with central fiber $X \xrightarrow{\varphi} Z$, where $\widetilde{X}$ is a first--order, locally trivial,
infinitesimal deformation of $X$.

An space classifying them is described in \cite{Hor74}. Let $\mathcal V=(V_i)$ be an open affine cover of $X$.
As usual, we let $C^{0}(\mathcal V,-)$ and $Z^{1}(\mathcal V,-)$ denote, respectively, the group of
$0$--cochains and $1$--cocycles with respect to the covering $\mathcal V$ and $\delta$ the coboundary map.
Let
\[
\varphi^{*}\Omega_{Z} \xrightarrow{\,D\varphi\,} \Omega_{X}
\]
be the homomorphism induced by $X \xrightarrow{\varphi} Z$. Set
\begin{equation}\label{2.2.1}
D(X,\varphi)
=
\frac{
\Bigl\{(g,\rho)\in
C^{0}\!\bigl(\mathcal V,\ \mathcal{H}om_{\mathcal O_X}(\varphi^{*}\Omega_{Z},\mathcal O_X)\bigr)
\times
Z^{1}\!\bigl(\mathcal V,\ \mathcal{H}om_{\mathcal O_X}(\Omega_{X},\mathcal O_X)\bigr)
\ \Big|\ 
\delta g = \rho\, D\varphi
\Bigr\}
}{
\Bigl\{(h\,D\varphi,\ \delta h)\ \Big|\ 
h\in C^{0}\!\bigl(\mathcal V,\ \mathcal{H}om_{\mathcal O_X}(\Omega_{X},\mathcal O_X)\bigr)
\Bigr\}
}.
\end{equation}

\begin{lemma}\cite[4.2]{Hor74}\label{D(X varphi)}
Let $X \xrightarrow{\varphi} Z$ be a morphism, where $X$ is a reduced connected scheme and $Z$ is a scheme.
Let $D(X,\varphi)$ be defined by \Cref{2.2.1}. Then
\begin{enumerate}
\item $D(X,\varphi)$ does not depend on the affine cover.
\item We have two exact sequences
\[
\operatorname{Hom}(\Omega_{X},\mathcal O_{X})
\xrightarrow{\,d\varphi\,}
\operatorname{Hom}(\varphi^{*}\Omega_{Z},\mathcal O_{X})
\longrightarrow
D(X,\varphi)
\longrightarrow
H^{1}\!\bigl(\mathcal{H}om(\Omega_{X},\mathcal O_{X})\bigr)
\xrightarrow{\,d\varphi\,}
H^{1}\!\bigl(\mathcal{H}om(\varphi^{*}\Omega_{Z},\mathcal O_{X})\bigr),
\]
\[
0 \longrightarrow
H^{1}\!\bigl(\mathcal{H}om(\Omega_{X/Z},\mathcal O_{X})\bigr)
\longrightarrow
D(X,\varphi)
\longrightarrow
H^{0}(N_{\varphi})
\longrightarrow
H^{2}\!\bigl(\mathcal{H}om(\Omega_{X/Z},\mathcal O_{X})\bigr),
\]
where $N_{\varphi}$ denotes the cokernel of
\[
\mathcal{H}om(\Omega_{X},\mathcal O_{X})
\xrightarrow{\,d\varphi\,}
\mathcal{H}om(\varphi^{*}\Omega_{Z},\mathcal O_{X}).
\]
Since is the kernel of $d\varphi$ is given by $\mathcal{H}om(\Omega_{X/Z},\mathcal O_{X})$, if the map
$\mathcal{H}om(\Omega_{X},\mathcal O_{X})
\xrightarrow{\,d\varphi\,}
\mathcal{H}om(\varphi^{*}\Omega_{Z},\mathcal O_{X})$
is injective then there is a natural isomorphism $D(X,\varphi)\simeq H^{0}(N_{\varphi})$. 
\end{enumerate}
\end{lemma}

\begin{proof}
    This lemma is derived directly from \Cref{D(X F)} where the morphism of coherent sheaves in question is \[
\mathcal{H}om(\Omega_{X},\mathcal O_{X})
\xrightarrow{\,d\varphi\,}
\mathcal{H}om(\varphi^{*}\Omega_{Z},\mathcal O_{X}).
\]
Note that the Kernel of this map is $\mathcal{H}om(\Omega_{X/Z},\mathcal O_{X})$ and the cokernel is $N_{\varphi}$ by definition.
\end{proof}


\begin{remark}\label{D(X_varphi) for nondegenerate varphi}
In \Cref{D(X varphi)}, If the map $\varphi$ is non-degenerate, meaning unramified on a dense open subset of $X$ (see \cite[Definition $3.4.5$]{Ser}) then $\mathcal{H}om(\Omega_{X/Z}, \mathcal{O}_X) = 0$ and consequently there is a natural isomorphism $D(X,\varphi)\simeq H^{0}(N_{\varphi})$.  
\end{remark}

\begin{remark}\label{dualizing_conormal_sequence}
Suppose that $Y$ is integral and $Z$ is smooth. Consider the exact sequence:
 
$$\mathcal{I}/\mathcal{I}^2 \to i^*\Omega_Z \xrightarrow{Di} \Omega_Y \to 0$$

Let $(\mathcal{I}/\mathcal{I}^2)'$ denote the image of the map 
$\mathcal{I}/\mathcal{I}^2 \to i^*\Omega_Z$. Then we have a short exact sequence

\begin{equation}\label{equisingular_conormal}
0 \to (\mathcal{I}/\mathcal{I}^2)' \to i^*\Omega_Z \xrightarrow{Di} \Omega_Y \to 0
\end{equation}
and a short exact sequence 

$$0 \to K \to \mathcal{I}/\mathcal{I}^2 \to (\mathcal{I}/\mathcal{I}^2)' \to 0 $$

Note that $K$ is a torsion sheaf. Now applying the functor $\mathcal{H}om(-,\mathcal{E})$ for a torsion free sheaf $\mathcal{E}$, we have an exact sequence 

$$0 \to \mathcal{H}om((\mathcal{I}/\mathcal{I}^2)', \mathcal{E}) \to \mathcal{H}om(\mathcal{I}/\mathcal{I}^2, \mathcal{E}) \to \mathcal{H}om(K, \mathcal{E}) = 0$$

Hence $\mathcal{H}om((\mathcal{I}/\mathcal{I}^2)', \mathcal{E}) \cong \mathcal{H}om(\mathcal{I}/\mathcal{I}^2, \mathcal{E})$

Applying the functor $\mathcal{H}om(-, \mathcal{E})$ for a torsion free sheaf $\mathcal{E}$, we get from \ref{equisingular_conormal}

\begin{align*}
0 \to \mathcal{H}om(\Omega_Y, \mathcal{E}) \xrightarrow{di}   \mathcal{H}om(i^*\Omega_Z, \mathcal{E}) \to \mathcal{H}om(\mathcal{I}/\mathcal{I}^2, \mathcal{E}) \to \mathcal{E}xt^1(\Omega_Y, \mathcal{E}) \to \mathcal{E}xt^1(i^*\Omega_Z, \mathcal{E}) = 0
\end{align*}

If $\mathcal{E} = \mathcal{O}_Y$, we get

\begin{align*}
0 \to \mathcal{H}om(\Omega_Y, \mathcal{O}_Y) \xrightarrow{di}   \mathcal{H}om(i^*\Omega_Z, \mathcal{O}_Y) \to \mathcal{H}om(\mathcal{I}/\mathcal{I}^2, \mathcal{O}_Y) \to \mathcal{E}xt^1(\Omega_Y, \mathcal{O}_Y) \to \mathcal{E}xt^1(i^*\Omega_Z, \mathcal{O}_Y) = 0
\end{align*}

Now in the notation of \cite[Proposition $1.1.9$, Theorem $1.1.10$]{Ser}, $T_Y = \mathcal{H}om(\Omega_Y, \mathcal{O}_Y)$, while $\mathcal{H}om(i^*\Omega_Z, \mathcal{O}_Y) = T_Z|_Y$. and 
$\mathcal{E}xt^1(\Omega_Y, \mathcal{O}_Y) = T_Y^1$. Hence comparing with \cite[Proposition $1.1.9$ (ii)]{Ser}, we have $\mathcal{H}om(\mathcal{I}/\mathcal{I}^2, \mathcal{O}_Y) = N_{Y/Z}$, which known as the normal sheaf of $Y$ in $Z$. The cokernel of $di$ (i.e, the image of the map $\mathcal{H}om(i^*\Omega_Z, \mathcal{O}_Y) \to \mathcal{H}om(\mathcal{I}/\mathcal{I}^2, \mathcal{O}_Y)$) is called the the equisingular normal sheaf and is denoted by $N_{Y/Z}'$. Therefore there are exact sequences

$$ 0 \to \mathcal{H}om(\Omega_Y, \mathcal{O}_Y) \xrightarrow{di}   \mathcal{H}om(i^*\Omega_Z, \mathcal{O}_Y) \to N_{Y/Z}' \to 0$$

There is an exact sequence 

$$0 \to N_{Y/Z}' \to N_{Y/Z} \to \mathcal{E}xt^1(\Omega_Y, \mathcal{O}_Y) \to 0$$

Now note that for a torsion free sheaf $\mathcal{E}$, $\operatorname{Hom}((\mathcal{I}/\mathcal{I}^2)', \mathcal{E}) = H^0(\mathcal{H}om((\mathcal{I}/\mathcal{I}^2)', \mathcal{E})) = H^0(\mathcal{H}om(\mathcal{I}/\mathcal{I}^2, \mathcal{E})) =  \operatorname{Hom}(\mathcal{I}/\mathcal{I}^2), \mathcal{E})$
Applying $\operatorname{Hom}(-, \mathcal{E})$ for a torsion free sheaf $\mathcal{E}$, to sequence \ref{equisingular_conormal}., we get 

\begin{align}\label{counting_reflexive_ribbons}
0 \to \operatorname{Hom}(\Omega_Y, \mathcal{E}) \xrightarrow{di}   \operatorname{Hom}(i^*\Omega_Z, \mathcal{E}) \to \operatorname{Hom}(\mathcal{I}/\mathcal{I}^2, \mathcal{E}) \to \operatorname{Ext}^1(\Omega_Y, \mathcal{E}) \to \operatorname{Ext}^1(i^*\Omega_Z, \mathcal{E}) 
\end{align}

\end{remark}



\begin{remark}\label{D(X_pi^*di)}

Consider the exact sequence 
$$0 \to (\mathcal{I}/\mathcal{I}^2)' \to i^*\Omega_Z \xrightarrow{Di} \Omega_Y \to 0$$

Applying $\pi^*$, and dualizing we get

$$0 \to \mathcal{H}om(\pi^*\Omega_Y, \mathcal{O}_X) \xrightarrow{\pi^*di} \mathcal{H}om(\varphi^*\Omega_Z, \mathcal{O}_X) \to \mathcal{H}om(\pi^*(\mathcal{I}/\mathcal{I}^2)', \mathcal{O}_X) \to \mathcal{E}xt^1(\pi^*\Omega_Y, \mathcal{O}_X) \to \mathcal{E}xt^1(\varphi^*\Omega_Z, \mathcal{O}_X) = 0 $$

The last term is zero since we are assuming $Z$ to be smooth. \par
Let the cokernel of $\pi^*di$ in the above sequence be $\mathcal{N}$ giving rise to two exact sequences

$$0 \to \mathcal{H}om(\pi^*\Omega_Y, \mathcal{O}_X) \xrightarrow{\pi^*di} \mathcal{H}om(\varphi^*\Omega_Z, \mathcal{O}_X) \to \mathcal{N} \to 0 $$

and 

\begin{equation}\label{N}
0 \to \mathcal{N} \to \mathcal{H}om(\pi^*(\mathcal{I}/\mathcal{I}^2)', \mathcal{O}_X) \to \mathcal{E}xt^1(\pi^*\Omega_Y, \mathcal{O}_X) \to 0 
\end{equation}
By \Cref{D(X varphi)}, we get a sheaf $D(X, \pi^*di)$ and there are two exact sequences
\[
\operatorname{Hom}(\pi^*\Omega_Y, \mathcal{O}_X))\longrightarrow \operatorname{Hom}(\varphi^*\Omega_Z, \mathcal{O}_X))\longrightarrow D(X,\pi^*di)
\longrightarrow H^{1}(\mathcal{H}om(\pi^*\Omega_Y, \mathcal{O}_X))\longrightarrow H^{1}(\mathcal{H}om(\varphi^*\Omega_Z, \mathcal{O}_X)),
\]
\[
0 \longrightarrow D(X,\pi^*di)\longrightarrow H^{0}(\mathcal N)
\longrightarrow 0.
\]
\end{remark}

\begin{remark}\label{D(X_pi_*di)}

Consider exact sequence 
$$0 \to (\mathcal{I}/\mathcal{I}^2)' \to i^*\Omega_Z \xrightarrow{Di} \Omega_Y \to 0$$

Applying the functor $\mathcal{H}om( -, \pi_*\mathcal{O}_X)$

$$0 \to \mathcal{H}om(\Omega_Y, \pi_*\mathcal{O}_X) \xrightarrow{(Di)'} \mathcal{H}om(i^*\Omega_Z, \pi_*\mathcal{O}_X) \to \mathcal{H}om(\mathcal{I}/\mathcal{I}^2, \pi_*\mathcal{O}_X) $$
Let the cokernel of $(Di)'$ in the above sequence be $\mathcal{M}$ giving a natural injection
$$0 \to H^0(\mathcal{M}) \to \operatorname{Hom}(\mathcal{I}/\mathcal{I}^2, \pi_*\mathcal{O}_X) $$
Note also that $\mathcal{M} = \pi_*\mathcal{N}$ since $\pi$ is a finite map and hence $\pi_*$ is exact.
By \Cref{D(X varphi)}, we get a sheaf $D(X, (Di)')$ and there are two exact sequences
\[
\operatorname{Hom}(\Omega_Y, \pi_*\mathcal{O}_X))\longrightarrow \operatorname{Hom}(i^*\Omega_Z, \pi_*\mathcal{O}_X))\longrightarrow D(X,(Di)')
\longrightarrow H^{1}(\mathcal{H}om(\Omega_Y, \pi_*\mathcal{O}_X))\longrightarrow H^{1}(\mathcal{H}om(i^*\Omega_Z, \pi_*\mathcal{O}_X)),
\]
\[
0 \longrightarrow D(X,(Di)')\longrightarrow H^{0}(\pi_*\mathcal N)
\longrightarrow 0.
\]
\end{remark}


% Lemma 3.4 (as in "Smoothing of ribbons over curves")
% Requires: \usepackage{tikz-cd}

\begin{remark}
    We have an exact sequence

    $$0 \to \mathcal{H}om(\Omega_{X},\mathcal O_{X})
\xrightarrow{\,d\pi\,}
\mathcal{H}om(\pi^{*}\Omega_{Y},\mathcal O_{X}) \to N_{\pi} \to 0$$

Similarly, we have an exact sequence 

$$0 \to \mathcal{H}om(\Omega_{X},\mathcal O_{X})
\xrightarrow{\,d\varphi\,}
\mathcal{H}om(\varphi^{*}\Omega_{Z},\mathcal O_{X}) \to N_{\varphi} \to 0$$Applying the snake lemma, we get an exact sequence 

$$ 0 \to N_{\pi} \to N_{\varphi} \to \mathcal{N} \to 0$$

Pushing it forward, we get 

$$ 0 \to \pi_*N_{\pi} \to \pi_*N_{\varphi} \to \pi_*\mathcal{N} \to 0$$

where by sequence \ref{N}, $\pi_*\mathcal{N}$ sits in an exact sequence as follows:

$$0 \to \mathcal{M} = \pi_*\mathcal{N} \to \mathcal{H}om(\mathcal{I}/\mathcal{I}^2, \pi_*\mathcal{O}_X) \to \pi_*\mathcal{E}xt^1(\pi^*\Omega_Y, \mathcal{O}_X) \to 0$$

$$0 \to \mathcal{M} = \pi_*\mathcal{N} \to \mathcal{H}om(\mathcal{I}/\mathcal{I}^2, \mathcal{O}_Y) \bigoplus \mathcal{H}om(\mathcal{I}/\mathcal{I}^2, \mathcal{E}) \to \pi_*\mathcal{E}xt^1(\pi^*\Omega_Y, \mathcal{O}_X) \to 0$$


If $\pi: X \to Y$ is an abelian cover so that $\pi_*\mathcal{O}_X = \mathcal{O}_Y \oplus_{\chi \in G^*-\{0\} }\mathcal{L}_{\chi}^{-1}$

$$0 \to \mathcal{M} = \pi_*\mathcal{N} \to \mathcal{H}om(\mathcal{I}/\mathcal{I}^2, \mathcal{O}_Y) \bigoplus \oplus_{\oplus_{\chi \in G^*-\{0\} } }\mathcal{H}om(\mathcal{I}/\mathcal{I}^2, \mathcal{L}_{\chi}^{-1}) \to \pi_*\mathcal{E}xt^1(\pi^*\Omega_Y, \mathcal{O}_X) \to 0$$

\end{remark}

% --- Derived adjunction + Ext/Tor discussion for a finite morphism ---
% Suggested bib entry:
% \bibitem[Stacks]{stacks-project}
% The Stacks Project Authors, \emph{Stacks Project}, \url{https://stacks.math.columbia.edu}.

\begin{lemma}[Derived adjunction for internal Hom {\cite[Tags 07A6, 07A4, 08J9]{stacks-project}}]
\label{lem:derived-adjunction-internal-hom}
Let $f : (X,\mathcal O_X)\to (Y,\mathcal O_Y)$ be a morphism of ringed topoi (in particular, a morphism
of schemes). For any $K\in D(\mathcal O_Y)$ and $L\in D(\mathcal O_X)$ there is a canonical isomorphism
in $D(\mathcal O_Y)$
\[
Rf_*\,R\mathcal H\!om_X\bigl(Lf^*K,\;L\bigr)\ \cong\ R\mathcal H\!om_Y\bigl(K,\;Rf_*L\bigr).
\]
\end{lemma}

\begin{proof}
Let $M\in D(\mathcal O_Y)$ be arbitrary. Using the adjunction $Lf^*\dashv Rf_*$ \cite[Tag 07A6]{stacks-project}
and the defining property of $R\mathcal H\!om$ \cite[Tag 08J9]{stacks-project}, we compute

\begin{align*}
\operatorname{Hom}_{D(\mathcal O_Y)}\!\bigl(M,\;Rf_*R\mathcal H\!om_X(Lf^*K,L)\bigr)
&\cong \operatorname{Hom}_{D(\mathcal O_X)}\!\bigl(Lf^*M,\;R\mathcal H\!om_X(Lf^*K,L)\bigr) \\
&\cong \operatorname{Hom}_{D(\mathcal O_X)}\!\bigl(Lf^*M\otimes_{\mathcal O_X}^{\mathbf L}Lf^*K,\;L\bigr) \\
&\cong \operatorname{Hom}_{D(\mathcal O_X)}\!\bigl(Lf^*(M\otimes_{\mathcal O_Y}^{\mathbf L}K),\;L\bigr)
\quad\text{by \cite[Tag 07A4]{stacks-project}} \\
&\cong \operatorname{Hom}_{D(\mathcal O_Y)}\!\bigl(M\otimes_{\mathcal O_Y}^{\mathbf L}K,\;Rf_*L\bigr)
\quad\text{by \cite[Tag 07A6]{stacks-project}} \\
&\cong \operatorname{Hom}_{D(\mathcal O_Y)}\!\bigl(M,\;R\mathcal H\!om_Y(K,Rf_*L)\bigr).
\end{align*}
again using the defining property of $R\mathcal H\!om$ \cite[Tag 08J9]{stacks-project}.
By Yoneda, this yields the claimed isomorphism in $D(\mathcal O_Y)$.
\end{proof}

\begin{remark}[Finite case and the Ext--Tor spectral sequence {\cite[Tag 0G1Z]{stacks-project}}]
\label{rem:finite-ext-tor}
Let $\pi:X\to Y$ be a finite morphism of schemes. For $\mathcal F\in D(\mathcal O_Y)$ and
$\mathcal G\in D(\mathcal O_X)$, Lemma~\ref{lem:derived-adjunction-internal-hom} gives
\[
R\pi_*\,R\mathcal H\!om_X\bigl(L\pi^*\mathcal F,\;\mathcal G\bigr)\ \cong\
R\mathcal H\!om_Y\bigl(\mathcal F,\;R\pi_*\mathcal G\bigr).
\]
Since $\pi$ is affine, $R\pi_*=\pi_*$ on complexes with quasi-coherent cohomology sheaves.
In particular, taking $\mathcal F=\Omega_Y$ and $\mathcal G=\mathcal O_X$ yields
\[
\pi_*\,R\mathcal H\!om_X\bigl(L\pi^*\Omega_Y,\;\mathcal O_X\bigr)
\ \cong\
R\mathcal H\!om_Y\bigl(\Omega_Y,\;\pi_*\mathcal O_X\bigr),
\]
hence on cohomology sheaves
\[
\pi_*\,\mathcal E\!xt_X^n\bigl(L\pi^*\Omega_Y,\mathcal O_X\bigr)
\ \cong\
\mathcal E\!xt_Y^n\bigl(\Omega_Y,\pi_*\mathcal O_X\bigr)\qquad(n\ge 0).
\]

Moreover, filter the object $L\pi^*\Omega_Y\simeq \Omega_Y\otimes_{\mathcal O_Y}^{\mathbf L}\mathcal O_X$
by truncations (the standard $t$--structure filtration). The associated graded pieces are the cohomology
sheaves of $L\pi^*\Omega_Y$, namely
\[
\mathcal H^{-q}(L\pi^*\Omega_Y)\ \cong\ \mathcal T\!or_q^{\mathcal O_Y}\bigl(\Omega_Y,\mathcal O_X\bigr)\qquad(q\ge 0).
\]
Applying the spectral sequence for $\operatorname{Ext}$ associated to a filtered complex
\cite[Tag 0G1Z]{stacks-project} (in the Grothendieck abelian category $\operatorname{Mod}(\mathcal O_X)$),
and pushing forward via the exact functor $\pi_*$, we obtain a spectral sequence
\[
E_2^{p,q}
=
\pi_*\,\mathcal E\!xt_X^{p}\!\Bigl(\mathcal T\!or_q^{\mathcal O_Y}(\Omega_Y,\mathcal O_X),\ \mathcal O_X\Bigr)
\ \Rightarrow\
\mathcal E\!xt_Y^{p+q}\bigl(\Omega_Y,\pi_*\mathcal O_X\bigr).
\]
In low degree this gives the exact edge sequence
\[
0\to
\pi_*\,\mathcal E\!xt_X^{1}\bigl(\pi^*\Omega_Y,\mathcal O_X\bigr)
\to
\mathcal E\!xt_Y^{1}\bigl(\Omega_Y,\pi_*\mathcal O_X\bigr)
\to
\pi_*\,\mathcal H\!om_X\!\Bigl(\mathcal T\!or^{\mathcal O_Y}_1(\Omega_Y,\mathcal O_X),\ \mathcal O_X\Bigr)\to\cdots,
\]
where $\pi^*\Omega_Y$ appears because $\mathcal H^0(L\pi^*\Omega_Y)\cong \pi^*\Omega_Y$.
\end{remark}


\begin{lemma}\label{lem:3.4}
In the conditions of \S3.1 there is a commutative diagram
\begin{equation}\label{eq:3.4.1}
\begin{tikzcd}[column sep=large,row sep=large]
D(X,\varphi) \arrow[r,"\sim"] \arrow[d,"\mu"'] &
H^{0}(N_{\varphi}) \arrow[d,"\phi"] & & \\
D(X,\pi_{*}d i) \arrow[r,"\sim"] \arrow[d,"\alpha'"', "\cong"] &
H^0(\mathcal{N}) \arrow[d,"\alpha", "\cong"'] \arrow[r, hook] & \operatorname{Hom}(\pi^*(\mathcal{I}/\mathcal{I}^2)', \mathcal{O}_X) \arrow[d, "\cong"] \arrow[r, hook] & \operatorname{Hom}(\pi^*(\mathcal{I}/\mathcal{I}^2), \mathcal{O}_X) \arrow[d, "\cong"] \\
D\bigl(Y,(D i)'\bigr) \arrow[r,"\sim"] &
H^0(\pi_*\mathcal{N}) \arrow[r, hook] & \operatorname{Hom}((\mathcal{I}/\mathcal{I}^2)', \pi_*\mathcal{O}_X) \arrow[r, hook] & \operatorname{Hom}((\mathcal{I}/\mathcal{I}^2), \pi_*\mathcal{O}_X) 
\end{tikzcd}
\end{equation}
where \(H^{0}(N_{\varphi}) \xrightarrow{\ \phi\ }
\operatorname{Hom}(\pi_{*}I/I^{2},\mathcal O_{X})\) is the map in global sections obtained from
\eqref{eq:3.3.1}, the composition
\(D(X,\varphi)\to \operatorname{Hom}(\pi_{*}I/I^{2},\mathcal O_{X})\) is
\([(g,\rho)] \mapsto g|_{\pi_{*}I/I^{2}}\), and if \(\alpha'\mu([(g,\rho)])=[(f,\varrho)]\)
then the composition
\(D(X,\varphi)\to \operatorname{Hom}(I/I^{2},\pi_{*}\mathcal O_{X})\) is
\([(g,\rho)] \mapsto f|_{I/I^{2}}\).
\end{lemma}

\begin{proof}
The spaces \(D(X,\varphi)\) and \(D(X,\pi_{*}d i)\) are associated, from Remark~2.4, respectively,
to the homomorphisms \(d\varphi\) and \(\pi_{*}d i\) in \eqref{eq:3.3.2}. The space
\(D\bigl(Y,(D i)'\bigr)\) is associated from Remark~2.4 to the homomorphism \((D i)'\) in the
exact sequence
\begin{equation}\label{eq:3.4.2}
0 \to \mathcal{H}om_{\mathcal O_Y}(\Omega_Y,\pi_{*}\mathcal O_X)
\xrightarrow{\ (D i)'\ }
\mathcal{H}om_{\mathcal O_Y}(i^{*}\Omega_Z,\pi_{*}\mathcal O_X)
\to
\mathcal{H}om_{\mathcal O_Y}(I/I^{2},\pi_{*}\mathcal O_X)
\to 0.
\end{equation}
So if \(\mathcal U=(U_i)\) is an open affine cover of \(Y\) and \(\mathcal V=(V_i=\pi^{-1}U_i)\) is
the induced open affine cover of \(X\), then \(D(X,\varphi)\) is given by \eqref{eq:2.2.1} and
\(D(X,\pi_{*}d i)\), \(D\bigl(Y,(D i)'\bigr)\) are obtained from \eqref{eq:2.4.1}.

The homomorphisms \(d\varphi\), \(\pi_{*}d i\) in \eqref{eq:3.3.2} and \((D i)'\) in \eqref{eq:3.4.2}
are injective. Therefore, from Remark~2.4, we have natural isomorphisms
\[
D(X,\varphi)\xrightarrow{\ \sim\ } H^{0}(N_{\varphi}),\qquad
D(X,\pi_{*}d i)\xrightarrow{\ \sim\ }\operatorname{Hom}(\pi_{*}I/I^{2},\mathcal O_X),\qquad
D\bigl(Y,(D i)'\bigr)\xrightarrow{\ \sim\ }\operatorname{Hom}(I/I^{2},\pi_{*}\mathcal O_X).
\]
From \eqref{eq:3.3.2} and the functorial character of \(D(X,-)\) in the exact sequences of
Remark~2.4 we see that there is a unique map \(D(X,\varphi)\xrightarrow{\ \mu\ }D(X,\pi_{*}d i)\)
making the upper square in \eqref{eq:3.4.1} commutative. Hence this map is
\[
[(g,\rho)] \ \longmapsto\ [(g,\rho\,D\pi)].
\]

From the adjunction isomorphism we obtain an isomorphism
\begin{equation}\label{eq:3.4.3}
D(X,\pi_{*}d i)\xrightarrow{\ \alpha'\ }\sim D\bigl(Y,(D i)'\bigr),
\qquad
[(g,\rho')] \longmapsto [(f,\varrho)].
\end{equation}
Moreover, the isomorphism \(D(X,\pi_{*}d i)\xrightarrow{\ \sim\ }
\operatorname{Hom}(\pi_{*}I/I^{2},\mathcal O_X)\) given by Remark~2.4 is defined by
\(
[(g,\rho')]\mapsto g|_{\pi_{*}I/I^{2}}
\),
i.e. for the restriction of the cochain \(g\) to \(\pi_{*}I/I^{2}\) we have
\(\delta(g|_{\pi_{*}I/I^{2}})=0\) and therefore there is a global section
\(g|_{\pi_{*}I/I^{2}}\in \operatorname{Hom}(\pi_{*}I/I^{2},\mathcal O_X)\).
Alike, the isomorphism \(D\bigl(Y,(D i)'\bigr)\xrightarrow{\ \sim\ }
\operatorname{Hom}(I/I^{2},\pi_{*}\mathcal O_X)\) is given by
\(
[(f,\varrho)]\mapsto f|_{I/I^{2}}.
\)

Now, since \(g\) corresponds to \(f\) by the adjunction isomorphism, we see that
\(g|_{\pi_{*}I/I^{2}}\) corresponds to \(f|_{I/I^{2}}\) by the isomorphism \(\alpha\) in \eqref{eq:3.4.1}.
So we have commutativity in the lower square of \eqref{eq:3.4.1} and the final assertions in the
Lemma follow as well.
\end{proof}

\section{Reflexive Ribbons}

\begin{definition}[Ropes with conormal sheaf]
\label{def:ribbon-conormal-sheaf}
Let \(Y\) be a reduced connected scheme, and let \(\mathcal E\) be a coherent
\(\mathcal O_Y\)-module. A \emph{ribbon over \(Y\)
with conormal sheaf \(\mathcal E\)} is a scheme \(\widetilde Y\), together with a closed
immersion
\[
    Y\hookrightarrow \widetilde Y,
\]
such that
\[
    \widetilde Y_{\mathrm{red}}=Y,
    \qquad
    \mathcal I_{Y/\widetilde Y}^{2}=0,
    \qquad
    \mathcal I_{Y/\widetilde Y}\simeq \mathcal E
\]
as \(\mathcal O_Y\)-modules. Equivalently, \(\widetilde Y\) fits into an exact sequence of sheaves of rings
\[
    0\longrightarrow \mathcal E
    \longrightarrow \mathcal O_{\widetilde Y}
    \longrightarrow \mathcal O_Y
    \longrightarrow 0,
\]
where \(\mathcal E\) is an ideal of square zero. If \(\mathcal E\) is reflexive of rank \(n-1\), then $\widetilde{Y}$ is called a \emph{reflexive \(n\)-rope} over \(Y\).
\end{definition}



\begin{remark}
\label{rem:properties-reflexive-rope}
Assume that \(Y\) is normal and that \(\mathcal E\) is reflexive.

\begin{enumerate}
    \item Since \(Y\) is normal, a reflexive sheaf on \(Y\) is locally free outside a closed
    subset of codimension at least two. Hence a reflexive rope is an ordinary rope on
    the open subset where \(\mathcal E\) is locally free. In particular, a reflexive rope
    is a genuine rope in codimension one.

    \item The scheme \(\widetilde Y\) is a square-zero thickening of \(Y\). More precisely,
    the closed immersion \(Y\hookrightarrow \widetilde Y\) is defined by the ideal
    \[
        \mathcal I_{Y/\widetilde Y}\simeq \mathcal E,
    \]
    and this ideal satisfies
    \[
        \mathcal I_{Y/\widetilde Y}^{2}=0.
    \]
    Thus the underlying topological space of \(\widetilde Y\) is the same as that of \(Y\),
    but the structure sheaf has an additional square-zero layer given by \(\mathcal E\).

    \item If \(\mathcal E\) has rank \(n-1\), then \(\widetilde Y\) has generic multiplicity
    \(n\) along \(Y\). Indeed, at the generic point \(\eta\) of \(Y\), one has
    \[
        \operatorname{length}_{\mathcal O_{Y,\eta}}
        \mathcal O_{\widetilde Y,\eta}
        =
        1+\operatorname{rank}\mathcal E
        =
        n.
    \]
    In particular, when \(\mathcal E\) has rank one, \(\widetilde Y\) is generically a
    double structure over \(Y\).

    \item \textcolor{blue}{Since \(Y\) is normal, \(\mathcal O_Y\) satisfies Serre's condition \(S_2\).
    Since \(\mathcal E\) is reflexive, \(\mathcal E\) also satisfies \(S_2\). From the exact
    sequence
    \[
        0\longrightarrow \mathcal E
        \longrightarrow \mathcal O_{\widetilde Y}
        \longrightarrow \mathcal O_Y
        \longrightarrow 0,
    \]
    the depth lemma implies that \(\mathcal O_{\widetilde Y}\) satisfies \(S_2\). Hence
    \(\widetilde Y\) is an \(S_2\)-scheme.}

    \item \textcolor{blue}{Reflexivity alone does not imply that \(\widetilde Y\) is Cohen--Macaulay in
    dimensions at least three. If \(Y\) is Cohen--Macaulay and \(\mathcal E\) is maximal
    Cohen--Macaulay as an \(\mathcal O_Y\)-module, then the same exact sequence shows
    that \(\widetilde Y\) is Cohen--Macaulay.}

   
\end{enumerate}
\end{remark}



\begin{proposition}[Injectivity of the conormal map for torsion-free square-zero ideals]
\label{prop:injectivity-conormal-map}
Let \(k\) be a field of characteristic \(0\). Let \(A\) be a reduced noetherian
\(k\)-algebra, and let
\[
    0\longrightarrow M\longrightarrow B\xrightarrow{q} A\longrightarrow 0
\]
be a square-zero extension of \(k\)-algebras; that is,
\[
    M=\ker(q)
    \qquad\text{and}\qquad
    M^2=0.
\]
Assume that \(M\) is torsion-free as an \(A\)-module, where torsion-free means that
multiplication by every nonzerodivisor of \(A\) is injective on \(M\). Then the natural
conormal map
\[
    M=M/M^2
    \longrightarrow
    \Omega_{B/k}\otimes_B A
\]
is injective.

Equivalently, we have an exact sequence the conormal sequence
\[
  0 \longrightarrow  M
    \longrightarrow
    \Omega_{B/k}\otimes_B A
    \longrightarrow
    \Omega_{A/k}
    \longrightarrow 0
\]
is left-exact.
\end{proposition}

\begin{proof}
Let
\[
    \delta:M\longrightarrow \Omega_{B/k}\otimes_B A
\]
be the conormal map. Explicitly, it is given by
\[
    m\longmapsto dm\otimes 1.
\]
Let
\[
    N:=\ker(\delta).
\]
We want to prove that \(N=0\).

We first reduce to the generic points of \(\operatorname{Spec}A\). Since \(A\) is noetherian,
it has only finitely many minimal primes; see \cite[Algebra, Lemma 10.31.6]{stacks-project}.
Let \(\mathfrak p\subset A\) be a minimal prime. Since \(A\) is reduced, the localization
\[
    A_{\mathfrak p}
\]
is a field; see \cite[Algebra, Lemma 10.25.1]{stacks-project}. We denote this field by
\[
    K:=A_{\mathfrak p}.
\]

Localizing the given square-zero extension at \(\mathfrak p\), we obtain an exact sequence
\[
    0\longrightarrow M_{\mathfrak p}
    \longrightarrow B_{\mathfrak p}
    \longrightarrow K
    \longrightarrow 0.
\]
The ideal \(M_{\mathfrak p}\) is still square-zero.

Because \(\operatorname{char}k=0\), every field extension of \(k\) is formally smooth over
\(k\); see \cite[Algebra, Lemma 10.158.8]{stacks-project}. Hence \(K/k\) is formally
smooth. Applying the infinitesimal lifting property for formal smoothness to the square-zero
surjection
\[
    B_{\mathfrak p}\twoheadrightarrow K
\]
and to the identity map \(K\to K\), we obtain a \(k\)-algebra section
\[
    K\longrightarrow B_{\mathfrak p}.
\]

Now apply the split conormal exact sequence for Kähler differentials to the surjection
\[
    B_{\mathfrak p}\twoheadrightarrow K
\]
with kernel \(M_{\mathfrak p}\). The split conormal sequence says that, since this surjection
has a \(k\)-algebra section, the sequence
\[
    0
    \longrightarrow
    M_{\mathfrak p}/M_{\mathfrak p}^{2}
    \longrightarrow
    \Omega_{B_{\mathfrak p}/k}\otimes_{B_{\mathfrak p}} K
    \longrightarrow
    \Omega_{K/k}
    \longrightarrow 0
\]
is exact; see \cite[Algebra, Lemma 10.131.10]{stacks-project}. Since
\[
    M_{\mathfrak p}^{2}=0,
\]
we have
\[
    M_{\mathfrak p}/M_{\mathfrak p}^{2}=M_{\mathfrak p}.
\]
Therefore the localized conormal map
\[
    M_{\mathfrak p}
    \longrightarrow
    \Omega_{B_{\mathfrak p}/k}\otimes_{B_{\mathfrak p}} K
\]
is injective.

Formation of Kähler differentials commutes with localization; see
\cite[Algebra, Lemma 10.131.8]{stacks-project}. Hence
\[
    \Omega_{B_{\mathfrak p}/k}
    \simeq
    \Omega_{B/k}\otimes_B B_{\mathfrak p}.
\]
After tensoring with \(K=A_{\mathfrak p}\), this identifies
\[
    \Omega_{B_{\mathfrak p}/k}\otimes_{B_{\mathfrak p}} K
\]
with
\[
    \left(\Omega_{B/k}\otimes_B A\right)_{\mathfrak p}.
\]
Thus the localized map
\[
    \delta_{\mathfrak p}:M_{\mathfrak p}
    \longrightarrow
    \left(\Omega_{B/k}\otimes_B A\right)_{\mathfrak p}
\]
is injective. Equivalently,
\[
    N_{\mathfrak p}=0
\]
for every minimal prime \(\mathfrak p\subset A\).

It remains to pass from generic injectivity to global injectivity. Let \(n\in N\). Since
\(N_{\mathfrak p}=0\) for every minimal prime \(\mathfrak p\), for each minimal prime
\(\mathfrak p\) there exists an element
\[
    s_{\mathfrak p}\in A\setminus \mathfrak p
\]
such that
\[
    s_{\mathfrak p}n=0.
\]
Equivalently, the annihilator ideal
\[
    \operatorname{Ann}_A(n)
\]
is not contained in any minimal prime of \(A\).

Since \(A\) has only finitely many minimal primes, prime avoidance
\cite[Algebra, Lemma 10.15.2]{stacks-project} implies that
\[
    \operatorname{Ann}_A(n)
\]
is not contained in the union of the minimal primes. Hence there exists an element
\[
    s\in \operatorname{Ann}_A(n)
\]
which is contained in no minimal prime of \(A\).

Because \(A\) is reduced, the set of zerodivisors of \(A\) is the union of its minimal
primes; see \cite[Algebra, Lemma 10.25.2]{stacks-project}. Therefore \(s\) is a
nonzerodivisor of \(A\). Since \(M\) is torsion-free and \(N\subset M\), multiplication by
\(s\) is injective on \(N\). But
\[
    sn=0
\]
by construction. Hence
\[
    n=0.
\]
Therefore every element of \(N\) is zero, so
\[
    N=0.
\]

Thus the conormal map
\[
    M\longrightarrow \Omega_{B/k}\otimes_B A
\]
is injective.
\end{proof}


\begin{proposition}[Classification of square-zero thickenings with torsion-free conormal sheaf]
\label{prop:classification-square-zero-thickenings}
Let \(k\) be a field of characteristic \(0\). Let \(Y\) be a reduced connected
noetherian \(k\)-scheme, and let \(\mathcal E\) be a coherent torsion-free
\(\mathcal O_Y\)-module. Consider pairs
\[
    (j:Y\hookrightarrow \widetilde Y,\alpha),
\]
where \(j\) is a closed immersion such that
\[
    \widetilde Y_{\mathrm{red}}=Y,
    \qquad
    \mathcal I_{Y/\widetilde Y}^{2}=0,
\]
and where
\[
    \alpha:\mathcal I_{Y/\widetilde Y}\xrightarrow{\sim}\mathcal E
\]
is an isomorphism of \(\mathcal O_Y\)-modules. Then the isomorphism classes of such pairs, where isomorphisms are required to be
over \(Y\) and compatible with the chosen identification with \(\mathcal E\), are
classified by
\[
    \operatorname{Ext}^{1}_{Y}(\Omega_{Y/k},\mathcal E).
\]

If the identification
\[
    \mathcal I_{Y/\widetilde Y}\simeq\mathcal E
\]
is not fixed, then the set of isomorphism classes is the quotient of
\[
    \operatorname{Ext}^{1}_{Y}(\Omega_{Y/k},\mathcal E)
\]
by the natural action of
\[
    \operatorname{Aut}_{\mathcal O_Y}(\mathcal E).
\]
\end{proposition}

\begin{proof}
Let
\[
    (j:Y\hookrightarrow \widetilde Y,\alpha)
\]
be such a square-zero thickening, and set
\[
    \mathcal K:=\mathcal I_{Y/\widetilde Y}.
\]
Thus
\[
    \mathcal K^{2}=0
\]
and
\[
    \alpha:\mathcal K\xrightarrow{\sim}\mathcal E.
\]

The closed immersion
\[
    j:Y\hookrightarrow \widetilde Y
\]
has the usual conormal sequence
\[
    \mathcal K/\mathcal K^{2}
    \longrightarrow
    \Omega_{\widetilde Y/k}\otimes_{\mathcal O_{\widetilde Y}}\mathcal O_Y
    \longrightarrow
    \Omega_{Y/k}
    \longrightarrow 0.
\]
Since \(\mathcal K^2=0\), this becomes
\[
    \mathcal K
    \longrightarrow
    \Omega_{\widetilde Y/k}|_Y
    \longrightarrow
    \Omega_{Y/k}
    \longrightarrow 0.
\]

We claim that the first map is injective. This is local on \(Y\). On an affine open
subset \(U=\operatorname{Spec}A\subset Y\), write
\[
    \widetilde U=\operatorname{Spec}B
\]
for the corresponding affine open subset of \(\widetilde Y\). Then we have a
square-zero extension
\[
    0\longrightarrow M\longrightarrow B\longrightarrow A\longrightarrow 0,
\]
where
\[
    M=\Gamma(U,\mathcal K).
\]
Since \(\mathcal K\simeq\mathcal E\), and since \(\mathcal E\) is torsion-free, the
\(A\)-module \(M\) is torsion-free. Also \(A\) is reduced because \(Y\) is reduced.
Therefore, by \Cref{prop:injectivity-conormal-map}, the conormal map
\[
    M
    \longrightarrow
    \Omega_{B/k}\otimes_B A
\]
is injective. Hence the sheaf map
\[
    \mathcal K
    \longrightarrow
    \Omega_{\widetilde Y/k}|_Y
\]
is injective.

Thus we obtain a short exact sequence
\[
    0
    \longrightarrow
    \mathcal K
    \longrightarrow
    \Omega_{\widetilde Y/k}|_Y
    \longrightarrow
    \Omega_{Y/k}
    \longrightarrow 0.
\]
Using the chosen isomorphism
\[
    \alpha:\mathcal K\xrightarrow{\sim}\mathcal E,
\]
this gives a short exact sequence
\[
    0
    \longrightarrow
    \mathcal E
    \longrightarrow
    \Omega_{\widetilde Y/k}|_Y
    \longrightarrow
    \Omega_{Y/k}
    \longrightarrow 0.
\]
Therefore the thickening determines an extension class
\[
    [e_{\widetilde Y}]
    \in
    \operatorname{Ext}^{1}_{Y}(\Omega_{Y/k},\mathcal E).
\]

Conversely, let
\[
    0
    \longrightarrow
    \mathcal E
    \longrightarrow
    \mathcal G
    \xrightarrow{p}
    \Omega_{Y/k}
    \longrightarrow 0
\]
be an extension of \(\Omega_{Y/k}\) by \(\mathcal E\). We construct from this a
square-zero thickening of \(Y\) with conormal sheaf \(\mathcal E\).

Let
\[
    d:\mathcal O_Y\longrightarrow \Omega_{Y/k}
\]
be the universal derivation. Define
\[
    \mathcal O_{\widetilde Y}
    :=
    \mathcal O_Y\times_{\Omega_{Y/k}}\mathcal G.
\]
Thus, for an open set \(U\subset Y\),
\[
    \mathcal O_{\widetilde Y}(U)
    =
    \left\{
        (f,g)\in \mathcal O_Y(U)\oplus\mathcal G(U)
        \;\middle|\;
        df=p(g)
    \right\}.
\]

We define multiplication on \(\mathcal O_{\widetilde Y}\) by
\[
    (f,g)\cdot(f',g')
    :=
    (ff',fg'+f'g).
\]
This multiplication is well-defined. Indeed, since
\[
    p(g)=df
    \qquad\text{and}\qquad
    p(g')=df',
\]
we have
\[
    p(fg'+f'g)
    =
    f\,p(g')+f'\,p(g)
    =
    f\,df'+f'\,df
    =
    d(ff').
\]
Thus
\[
    (ff',fg'+f'g)
\]
again lies in the fiber product
\[
    \mathcal O_Y\times_{\Omega_{Y/k}}\mathcal G.
\]
With this multiplication, \(\mathcal O_{\widetilde Y}\) is a sheaf of rings whose unit is
\[
    (1,0).
\]

The projection onto the first factor gives a surjection
\[
    \mathcal O_{\widetilde Y}
    \longrightarrow
    \mathcal O_Y,
    \qquad
    (f,g)\longmapsto f.
\]
Its kernel consists of pairs
\[
    (0,g)
\]
with
\[
    p(g)=0.
\]
Since
\[
    \ker(p)=\mathcal E,
\]
we get
\[
    \ker(\mathcal O_{\widetilde Y}\to\mathcal O_Y)\simeq\mathcal E.
\]
Moreover, if \(e,e'\) are local sections of \(\mathcal E\), then
\[
    (0,e)\cdot(0,e')=(0,0).
\]
Hence this kernel is an ideal of square zero.

Therefore the ringed space defined by \(\mathcal O_{\widetilde Y}\) is a square-zero
thickening
\[
    Y\hookrightarrow \widetilde Y
\]
whose ideal sheaf is identified with \(\mathcal E\).

These two constructions give the desired classification of marked square-zero thickenings
by
\[
    \operatorname{Ext}^{1}_{Y}(\Omega_{Y/k},\mathcal E).
\]

Finally, suppose the identification of the square-zero ideal with \(\mathcal E\) is not fixed.
Changing an identification
\[
    \mathcal I_{Y/\widetilde Y}\xrightarrow{\sim}\mathcal E
\]
by an automorphism
\[
    \sigma\in \operatorname{Aut}_{\mathcal O_Y}(\mathcal E)
\]
pushes out the corresponding extension class by \(\sigma\). Thus
\[
    \operatorname{Aut}_{\mathcal O_Y}(\mathcal E)
\]
acts naturally on
\[
    \operatorname{Ext}^{1}_{Y}(\Omega_{Y/k},\mathcal E),
\]
and the unmarked isomorphism classes are obtained by taking the quotient by this action.
\end{proof}



\begin{proposition}[Morphisms from square-zero thickenings to an ambient scheme]
\label{prop:square-zero-thickenings-and-morphisms}
Let
\[
    i:Y\hookrightarrow Z
\]
be a closed immersion of schemes, and let
\[
    I=I_{Y/Z}\subset \mathcal O_Z
\]
be its ideal sheaf. Let \(\mathcal E\) be an \(\mathcal O_Y\)-module. Consider triples
\[
    (\widetilde Y,j,\alpha),
\]
where
\[
    j:Y\hookrightarrow \widetilde Y
\]
is a closed immersion defined by a square-zero ideal
\[
    \mathcal K:=\mathcal I_{Y/\widetilde Y},
    \qquad
    \mathcal K^2=0,
\]
and
\[
    \alpha:\mathcal K\xrightarrow{\sim}\mathcal E
\]
is an isomorphism of \(\mathcal O_Y\)-modules.

Then the following hold.

\begin{enumerate}
    \item Isomorphism classes of pairs
    \[
        \left((\widetilde Y,j,\alpha),\widetilde i\right),
    \]
    where
    \[
        \widetilde i:\widetilde Y\to Z
    \]
    is a morphism satisfying
    \[
        \widetilde i\circ j=i,
    \]
    are naturally in bijection with homomorphisms
    \[
        \tau:I/I^2\longrightarrow \mathcal E.
    \]

    \item Under this correspondence, the scheme-theoretic image
    \[
        W:=\operatorname{Im}(\widetilde i)\subset Z
    \]
    has ideal sheaf
    \[
        I_{W/Z}
        =
        \ker\left(
            I
            \longrightarrow
            I/I^2
            \xrightarrow{\tau}
            \mathcal E
        \right).
    \]
    Consequently,
    \[
        I^2\subset I_{W/Z}\subset I,
    \]
    and the ideal of \(Y\) inside \(W\) is
    \[
        \mathcal I_{Y/W}
        \simeq
        \operatorname{im}(\tau).
    \]

    \item The morphism
    \[
        \widetilde i:\widetilde Y\to Z
    \]
    is a closed immersion if and only if
    \[
        \tau:I/I^2\to \mathcal E
    \]
    is surjective.

    In particular, embedded square-zero thickenings
    \[
        \widetilde Y\subset Z
    \]
    with
    \[
        \mathcal I_{Y/\widetilde Y}\simeq\mathcal E
    \]
    correspond precisely to surjections
    \[
        I/I^2\twoheadrightarrow \mathcal E.
    \]
\end{enumerate}
\end{proposition}

\begin{proof}
We first explain how a morphism
\[
    \widetilde i:\widetilde Y\to Z
\]
extending \(i:Y\hookrightarrow Z\) produces a homomorphism
\[
    \tau:I/I^2\to \mathcal E.
\]

Let
\[
    \mathcal K=\mathcal I_{Y/\widetilde Y}.
\]
Since \(\widetilde i\circ j=i\), the morphism on structure sheaves
\[
    \widetilde i^\#:\mathcal O_Z\longrightarrow \widetilde i_*\mathcal O_{\widetilde Y}
\]
reduces modulo \(\mathcal K\) to the map
\[
    i^\#:\mathcal O_Z\longrightarrow i_*\mathcal O_Y.
\]
Therefore, if \(a\) is a local section of \(I\), then \(i^\#(a)=0\), and hence
\[
    \widetilde i^\#(a)\in \mathcal K.
\]
Using the fixed isomorphism
\[
    \alpha:\mathcal K\xrightarrow{\sim}\mathcal E,
\]
we define
\[
    \lambda_{\widetilde i}(a)
    :=
    \alpha\bigl(\widetilde i^\#(a)\bigr)
    \in \mathcal E.
\]

We claim that \(\lambda_{\widetilde i}\) kills \(I^2\). Indeed, if \(a,b\) are local
sections of \(I\), then
\[
    \widetilde i^\#(a),\widetilde i^\#(b)\in \mathcal K.
\]
Since \(\mathcal K^2=0\), it follows that
\[
    \widetilde i^\#(ab)
    =
    \widetilde i^\#(a)\widetilde i^\#(b)
    =
    0.
\]
Thus
\[
    \lambda_{\widetilde i}(ab)=0.
\]
Therefore \(\lambda_{\widetilde i}\) descends to a homomorphism
\[
    \tau_{\widetilde i}:I/I^2\longrightarrow \mathcal E.
\]

This homomorphism is \(\mathcal O_Y\)-linear. To see this, let \(f\) be a local section
of \(\mathcal O_Z\), and let \(a\) be a local section of \(I\). Then
\[
    \widetilde i^\#(fa)
    =
    \widetilde i^\#(f)\widetilde i^\#(a).
\]
The section \(\widetilde i^\#(a)\) lies in \(\mathcal K\). Since \(\mathcal K^2=0\), the
\(\mathcal O_{\widetilde Y}\)-module structure on \(\mathcal K\) factors through
\[
    \mathcal O_{\widetilde Y}/\mathcal K\simeq\mathcal O_Y.
\]
Thus
\[
    \widetilde i^\#(f)\widetilde i^\#(a)
    =
    \overline f\cdot \widetilde i^\#(a),
\]
where \(\overline f\) is the image of \(f\) in \(\mathcal O_Y\). Applying \(\alpha\), we get
\[
    \tau_{\widetilde i}(\overline{fa})
    =
    \overline f\cdot \tau_{\widetilde i}(\overline a).
\]
Hence
\[
    \tau_{\widetilde i}:I/I^2\to\mathcal E
\]
is \(\mathcal O_Y\)-linear.

We now construct the inverse correspondence.

Let
\[
    \tau:I/I^2\to\mathcal E
\]
be an \(\mathcal O_Y\)-module homomorphism. Let
\[
    Y^{(1)}_Z:=V(I^2)\subset Z
\]
be the first infinitesimal neighborhood of \(Y\) in \(Z\). Its structure sheaf fits into
the square-zero extension
\[
    0\longrightarrow I/I^2
    \longrightarrow \mathcal O_{Y^{(1)}_Z}
    \longrightarrow \mathcal O_Y
    \longrightarrow 0.
\]
We push out this square-zero extension along
\[
    \tau:I/I^2\to\mathcal E.
\]

Concretely, define a sheaf of rings \(\mathcal O_{\widetilde Y_\tau}\) on the underlying
topological space of \(Y\) by
\[
    \mathcal O_{\widetilde Y_\tau}
    :=
    \left(
        \mathcal O_{Y^{(1)}_Z}\oplus \mathcal E
    \right)
    \big/
    \left\langle
        (u,-\tau(u))\mid u\in I/I^2
    \right\rangle.
\]
Here \(I/I^2\) is viewed as a square-zero ideal in \(\mathcal O_{Y^{(1)}_Z}\), and
\(\mathcal E\) is placed as a square-zero ideal. The multiplication is the one induced
from the standard square-zero multiplication:
\[
    (a,e)(a',e')
    =
    \bigl(
        aa',
        \overline a\, e'
        +
        \overline{a'}\, e
    \bigr),
\]
where \(\overline a,\overline{a'}\) denote the images in \(\mathcal O_Y\).

By construction, there is an exact sequence of sheaves of rings and modules
\[
    0\longrightarrow \mathcal E
    \longrightarrow \mathcal O_{\widetilde Y_\tau}
    \longrightarrow \mathcal O_Y
    \longrightarrow 0,
\]
and \(\mathcal E\) is an ideal of square zero. Therefore
\[
    j_\tau:Y\hookrightarrow \widetilde Y_\tau
\]
is a square-zero thickening with conormal sheaf \(\mathcal E\).

The natural quotient map
\[
    \mathcal O_Z
    \longrightarrow
    \mathcal O_Z/I^2
    =
    \mathcal O_{Y^{(1)}_Z}
\]
followed by
\[
    \mathcal O_{Y^{(1)}_Z}\to \mathcal O_{\widetilde Y_\tau}
\]
defines a morphism
\[
    \widetilde i_\tau:\widetilde Y_\tau\to Z.
\]
Modulo the square-zero ideal \(\mathcal E\), this morphism restricts to the original
closed immersion
\[
    i:Y\hookrightarrow Z.
\]
Hence
\[
    \widetilde i_\tau\circ j_\tau=i.
\]

We now check that the two constructions are inverse to each other.

Starting from a homomorphism \(\tau:I/I^2\to\mathcal E\), the construction above gives
\[
    \widetilde i_\tau:\widetilde Y_\tau\to Z.
\]
If \(a\) is a local section of \(I\), then its image in
\(\mathcal O_{Y^{(1)}_Z}\) is the class \(\overline a\in I/I^2\). In the pushout
\(\mathcal O_{\widetilde Y_\tau}\), this class becomes \(\tau(\overline a)\in\mathcal E\).
Therefore the homomorphism associated to \(\widetilde i_\tau\) is exactly \(\tau\).

Conversely, suppose we start with a pair
\[
    \left((\widetilde Y,j,\alpha),\widetilde i\right).
\]
The morphism
\[
    \widetilde i^\#:\mathcal O_Z\to\widetilde i_*\mathcal O_{\widetilde Y}
\]
kills \(I^2\). Hence it
factors through
\[
    \mathcal O_{Y^{(1)}_Z}=\mathcal O_Z/I^2.
\]
We obtain a morphism of square-zero extensions
\[
\begin{tikzcd}
    0 \arrow[r] &
    I/I^2 \arrow[r] \arrow[d,"\tau_{\widetilde i}"] &
    \mathcal O_{Y^{(1)}_Z} \arrow[r] \arrow[d] &
    \mathcal O_Y \arrow[r] \arrow[d,equal] &
    0
    \\
    0 \arrow[r] &
    \mathcal E \arrow[r] &
    \mathcal O_{\widetilde Y} \arrow[r] &
    \mathcal O_Y \arrow[r] &
    0 .
\end{tikzcd}
\]
By the universal property of the pushout, this induces a morphism
\[
    \mathcal O_{\widetilde Y_{\tau_{\widetilde i}}}
    \longrightarrow
    \mathcal O_{\widetilde Y}.
\]
This morphism is the identity on the quotient \(\mathcal O_Y\) and the identity on the
square-zero ideal \(\mathcal E\). Hence it is an isomorphism of sheaves of rings. Therefore
the original pair is recovered from \(\tau_{\widetilde i}\). This proves the bijection in
part \((1)\).

Next we compute the ideal of the scheme-theoretic image. Let
\[
    W=\operatorname{Im}(\widetilde i_\tau)\subset Z.
\]
By definition,
\[
    I_{W/Z}
    =
    \ker\left(
        \mathcal O_Z\to \mathcal O_{\widetilde Y_\tau}
    \right).
\]
If a local section \(f\in\mathcal O_Z\) maps to zero in \(\mathcal O_{\widetilde Y_\tau}\),
then its image in \(\mathcal O_Y\) is zero, so \(f\in I\). For \(a\in I\), its image in
\(\mathcal O_{\widetilde Y_\tau}\) is precisely
\[
    \tau(\overline a)\in\mathcal E,
\]
where \(\overline a\) denotes the class of \(a\) in \(I/I^2\). Therefore
\[
    a\in I_{W/Z}
    \quad\Longleftrightarrow\quad
    \tau(\overline a)=0.
\]
Thus
\[
    I_{W/Z}
    =
    \ker\left(
        I
        \longrightarrow
        I/I^2
        \xrightarrow{\tau}
        \mathcal E
    \right).
\]

Since this ideal contains \(I^2\) and is contained in \(I\), we have
\[
    I^2\subset I_{W/Z}\subset I.
\]
Moreover, the ideal of \(Y\) in \(W\) is
\[
    \mathcal I_{Y/W}=I/I_{W/Z}.
\]
By the first isomorphism theorem,
\[
    I/I_{W/Z}\simeq\operatorname{im}(\tau).
\]
Hence
\[
    \mathcal I_{Y/W}\simeq \operatorname{im}(\tau).
\]

Finally, we prove the closed immersion criterion. The morphism
\[
    \widetilde i_\tau:\widetilde Y_\tau\to Z
\]
is a closed immersion if and only if the corresponding morphism of sheaves of rings
\[
    \mathcal O_Z\to \mathcal O_{\widetilde Y_\tau}
\]
is surjective. Since
\[
    \mathcal O_Z\to\mathcal O_Y
\]
is already surjective, the only possible failure of surjectivity lies in the square-zero
ideal \(\mathcal E\subset\mathcal O_{\widetilde Y_\tau}\). The image of \(I\) inside
\(\mathcal E\) is exactly
\[
    \operatorname{im}(\tau).
\]
Therefore
\[
    \mathcal O_Z\to \mathcal O_{\widetilde Y_\tau}
\]
is surjective if and only if
\[
    \operatorname{im}(\tau)=\mathcal E,
\]
i.e. if and only if
\[
    \tau:I/I^2\to\mathcal E
\]
is surjective.

Thus \(\widetilde i_\tau\) is a closed immersion if and only if \(\tau\) is surjective.
This proves the proposition.
\end{proof}


\section{The map connecting deformations and reflexive ribbons}

\begin{lemma}[The induced map to the two conormal Hom groups]
\label{lem:T1-map-to-conormal-Hom}
Let \(k\) be an algebraically closed field, and let
\[
    X \xrightarrow{\pi} Y \xrightarrow{i} Z
\]
be morphisms of \(k\)-schemes, with \(\pi\) finite and \(i\) a closed immersion.
Let \(I=I_{Y/Z}\) be the ideal sheaf of \(Y\) in \(Z\), and set
\[
    \varphi=i\circ \pi.
\]
Assume that the finite \(\mathcal O_Y\)-algebra \(\pi_*\mathcal O_X\) admits a splitting
\[
    \pi_*\mathcal O_X \simeq \mathcal O_Y\oplus \mathcal E
\]
as an \(\mathcal O_Y\)-module. Then there is a natural homomorphism
\[
    \Theta=\Theta_1\oplus \Theta_2:
    T^1(X/Z)
    \longrightarrow
    \operatorname{Hom}_{\mathcal O_Y}(I/I^2,\mathcal O_Y)
    \oplus
    \operatorname{Hom}_{\mathcal O_Y}(I/I^2,\mathcal E),
\]
where
\[
    T^1(X/Z):=\operatorname{Ext}^1_X(L_{X/Z},\mathcal O_X).
\]
Equivalently, \(\Theta\) is the composite
\[
\begin{tikzcd}[column sep=large]
    T^1(X/Z)
    \arrow[r]
    &
    \operatorname{Ext}^1_X(L\pi^*L_{Y/Z},\mathcal O_X)
    \arrow[r,"\sim"]
    &
    \operatorname{Hom}_{\mathcal O_Y}(I/I^2,\pi_*\mathcal O_X)
    \arrow[r,"\sim"]
    &
    \operatorname{Hom}(I/I^2,\mathcal O_Y)
    \oplus
    \operatorname{Hom}(I/I^2,\mathcal E).
\end{tikzcd}
\]
\end{lemma}

\begin{proof}
We use the cotangent complex. By Wehler's convention for tangent cohomology,
\[
    T^i(X/Z,\mathcal M)
    =
    \operatorname{Ext}^i_X(L_{X/Z},\mathcal M),
\]
and in particular
\[
    T^1(X/Z)=\operatorname{Ext}^1_X(L_{X/Z},\mathcal O_X)
\]
is the tangent space to the functor of first-order deformations of \(X\) over the
fixed target \(Z\); see \cite[\S 1.1--1.2]{Wehler1986}.

Apply the transitivity triangle for the composable morphisms
\[
    X \xrightarrow{\pi} Y \xrightarrow{i} Z.
\]
It gives a distinguished triangle
\[
    L\pi^*L_{Y/Z}
    \longrightarrow
    L_{X/Z}
    \longrightarrow
    L_{X/Y}
    \xrightarrow{+1}.
\]
Equivalently, in Wehler's notation, every morphism \(g:X\to Y\) over \(Z\)
induces a short exact sequence in the derived category, hence long exact
sequences in tangent cohomology; see \cite[\S 1.3(e)]{Wehler1986}. Applying
\(\mathbf R\mathcal Hom_X(-,\mathcal O_X)\) to the first arrow of the triangle gives
a canonical homomorphism
\[
    \operatorname{Ext}^1_X(L_{X/Z},\mathcal O_X)
    \longrightarrow
    \operatorname{Ext}^1_X(L\pi^*L_{Y/Z},\mathcal O_X).
\]
Thus we have a natural map
\[
    T^1(X/Z)
    \longrightarrow
    \operatorname{Ext}^1_X(L\pi^*L_{Y/Z},\mathcal O_X).
\]

\textcolor{blue}{Since \(i:Y\hookrightarrow Z\) is a closed immersion with ideal \(I\), the
relative cotangent complex is
\[
    L_{Y/Z}\simeq (I/I^2)[1];
\]
see, for instance, \cite[Ch.~III, Proposition~2.2.4]{Illusie1971} or
\cite[Proposition~3.2.9]{Sernesi2006}.}
We do not use the stronger assertion $$L_{Y/Z}\simeq (I/I^2)[1],$$ which would require the closed immersion $i:Y\hookrightarrow Z$ to be a local complete intersection immersion. Instead, we use only the first truncation of the cotangent complex. Locally on $Z$, the closed immersion $i:Y\hookrightarrow Z$ is given by a surjection of rings $A\to B=A/I$. For a surjective ring map $A\to B$ with kernel $I$, the cotangent complex satisfies $$H^0(L_{B/A})=0$$ and $$H^{-1}(L_{B/A})\simeq I/I^2;$$ see \cite[The Cotangent Complex, Lemma~92.11.2]{stacks-project}. Hence, after sheafification, one has $$H^0(L_{Y/Z})=0$$ and $$H^{-1}(L_{Y/Z})\simeq I/I^2.$$ By the canonical truncation triangle in the derived category $D(\mathcal O_Y)$, one has
$$
\tau_{\leq -2}L_{Y/Z}
\longrightarrow
L_{Y/Z}
\longrightarrow
\tau_{\geq -1}L_{Y/Z}
\xrightarrow{+1};
$$
see \cite[Derived Categories, Remark~13.12.4]{stacks-project}. Since $\tau_{\geq -1}L_{Y/Z}$ has only one nonzero cohomology sheaf, namely $I/I^2$ in degree $-1$, it follows that
$$
\tau_{\geq -1}L_{Y/Z}\simeq (I/I^2)[1].
$$
Set $$K:=\tau_{\leq -2}L_{Y/Z}.$$ Then $K\in D^{\leq -2}(\mathcal O_Y)$, while any $\mathcal O_Y$-module $\mathcal M$ is regarded as an object of $D^{\geq 0}(\mathcal O_Y)$. Therefore, by the standard Ext-vanishing for objects with bounded cohomological degrees, one has
$$
\operatorname{Ext}^n_Y(K,\mathcal M)=0
\qquad\text{for } n<2;
$$
see \cite[Derived Categories, Lemma~13.27.3]{stacks-project}. Applying the functor $\operatorname{Ext}^{\bullet}_Y(-,\mathcal M)$ to the truncation triangle gives an isomorphism
$$
\operatorname{Ext}^1_Y(L_{Y/Z},\mathcal M)
\simeq
\operatorname{Ext}^1_Y((I/I^2)[1],\mathcal M).
$$
Finally, using the shift convention for Ext groups, one obtains
$$
\operatorname{Ext}^1_Y((I/I^2)[1],\mathcal M)
=
\operatorname{Ext}^0_Y(I/I^2,\mathcal M)
=
\operatorname{Hom}_{\mathcal O_Y}(I/I^2,\mathcal M).
$$
Thus, for every $\mathcal O_Y$-module $\mathcal M$, and in particular for $\mathcal M=\pi_*\mathcal O_X$ or for a direct summand $\mathcal M=\mathcal E$, one has
$$
T^1(Y/Z,\mathcal M)
=
\operatorname{Ext}^1_Y(L_{Y/Z},\mathcal M)
\simeq
\operatorname{Hom}_{\mathcal O_Y}(I/I^2,\mathcal M).
$$

Therefore
\[
\begin{aligned}
    \operatorname{Ext}^1_X(L\pi^*L_{Y/Z},\mathcal O_X)
    &\simeq
    \operatorname{Ext}^1_X(L\pi^*((I/I^2)[1]),\mathcal O_X)  \\
    &\simeq
    \operatorname{Hom}_{D(X)}(L\pi^*(I/I^2),\mathcal O_X).
\end{aligned}
\]
By derived adjunction,
\[
    \operatorname{Hom}_{D(X)}(L\pi^*(I/I^2),\mathcal O_X)
    \simeq
    \operatorname{Hom}_{D(Y)}(I/I^2,R\pi_*\mathcal O_X).
\]
Since \(\pi\) is finite, \(R\pi_*\mathcal O_X=\pi_*\mathcal O_X\). Hence
\[
    \operatorname{Ext}^1_X(L\pi^*L_{Y/Z},\mathcal O_X)
    \simeq
    \operatorname{Hom}_{\mathcal O_Y}(I/I^2,\pi_*\mathcal O_X).
\]

Finally, using the given splitting
\[
    \pi_*\mathcal O_X\simeq \mathcal O_Y\oplus \mathcal E,
\]
we obtain
\[
    \operatorname{Hom}_{\mathcal O_Y}(I/I^2,\pi_*\mathcal O_X)
    \simeq
    \operatorname{Hom}_{\mathcal O_Y}(I/I^2,\mathcal O_Y)
    \oplus
    \operatorname{Hom}_{\mathcal O_Y}(I/I^2,\mathcal E).
\]
The composite is the required homomorphism \(\Theta=\Theta_1\oplus\Theta_2\).
\end{proof}

\begin{remark}
When \(X\) and \(Y\) are smooth and the deformation of \(\varphi=i\circ\pi\) is
locally trivial, \(T^1(X/Z)\) is represented by the usual global normal space
\(H^0(N_\varphi)\). In this case the map \(\Theta\) recovers the homomorphism
\[
    H^0(N_\varphi)
    \longrightarrow
    \operatorname{Hom}(\pi^*(I/I^2),\mathcal O_X)
    \simeq
    \operatorname{Hom}(I/I^2,\pi_*\mathcal O_X)
\]
used by Gonz\'alez. After the trace decomposition
\(\pi_*\mathcal O_X=\mathcal O_Y\oplus\mathcal E\), this gives the two components
\[
    H^0(N_\varphi)\to \operatorname{Hom}(I/I^2,\mathcal O_Y),
    \qquad
    H^0(N_\varphi)\to \operatorname{Hom}(I/I^2,\mathcal E),
\]
as in \cite[Proposition~3.7]{Gonzalez2005}; see also
\cite[Proposition~1.2]{GallegoGonzalezPurnaprajna2010}.
\end{remark}

\begin{theorem}[Central fiber of the image of a first-order deformation]
\label{thm:central-fiber-image-singular-source}
Let \(k\) be an algebraically closed field and let
\[
    X \xrightarrow{\pi} Y \xrightarrow{i} Z
\]
be morphisms of \(k\)-schemes, where \(\pi\) is finite and \(i\) is a closed immersion.
Let
\[
    \varphi=i\circ \pi:X\to Z,
\]
and let \(I=I_{Y/Z}\subset \mathcal O_Z\) be the ideal sheaf of \(Y\) in \(Z\).
Assume that
\[
    \pi_*\mathcal O_X\simeq \mathcal O_Y\oplus \mathcal E
\]
as an \(\mathcal O_Y\)-module.

Let
\[
    \Delta=\operatorname{Spec} k[\epsilon]/(\epsilon^2),
\]
and let
\[
    \widetilde\varphi:\widetilde X\longrightarrow Z\times \Delta
\]
be a first-order deformation of \(\varphi\) over the fixed target \(Z\), i.e. assume that
\(\widetilde X\) is flat over \(\Delta\), its central fiber is \(X\), and the restriction of
\(\widetilde\varphi\) to the central fiber is \(\varphi\).

Let
\[
    \widetilde Y:=\operatorname{Im}(\widetilde\varphi)\subset Z\times\Delta
\]
be the scheme-theoretic image of \(\widetilde\varphi\), and let
\[
    \widetilde Y_0:=\widetilde Y\times_{\Delta}\operatorname{Spec}k
    \subset Z
\]
be its central fiber.

Associated to \(\widetilde\varphi\) there is a natural homomorphism
\[
    \nu_{\widetilde\varphi}: I/I^2\longrightarrow \pi_*\mathcal O_X.
\]
Using the splitting
\[
    \pi_*\mathcal O_X\simeq \mathcal O_Y\oplus \mathcal E,
\]
write
\[
    \nu_{\widetilde\varphi}
    =
    \nu_1\oplus \nu_2,
\]
where
\[
    \nu_1:I/I^2\to \mathcal O_Y,
    \qquad
    \nu_2:I/I^2\to \mathcal E.
\]
Then the ideal sheaf of \(\widetilde Y_0\) in \(Z\) is
\[
    I_{\widetilde Y_0/Z}
    =
    \ker\left(
        I
        \longrightarrow
        I/I^2
        \xrightarrow{\nu_2}
        \mathcal E
    \right).
\]
In particular,
\[
    I^2\subset I_{\widetilde Y_0/Z}\subset I,
\]
and hence
\[
    Y\subset \widetilde Y_0\subset Y^{(1)}_Z,
\]
where \(Y^{(1)}_Z\) denotes the first infinitesimal neighborhood of \(Y\) in \(Z\).
\end{theorem}

\begin{proof}
The proof is local on \(Z\), so let
\[
    W=\operatorname{Spec}A\subset Z
\]
be an affine open subset and set
\[
    U=W\cap Y=\operatorname{Spec}(A/I),
\]
where, by abuse of notation, \(I\subset A\) denotes the ideal defining \(U\subset W\).
Let
\[
    V=\pi^{-1}(U)=\operatorname{Spec}B.
\]
Since \(\pi\) is finite, we have
\[
    B=\Gamma(U,\pi_*\mathcal O_X).
\]
Using the splitting \(\pi_*\mathcal O_X\simeq \mathcal O_Y\oplus \mathcal E\), after restricting to
\(U\) we may write
\[
    B\simeq A/I\oplus M,
    \qquad
    M=\Gamma(U,\mathcal E).
\]

Let
\[
    R=k[\epsilon]/(\epsilon^2).
\]
Then
\[
    W\times\Delta=\operatorname{Spec}(A\otimes_k R)
    =
    \operatorname{Spec}(A\oplus A\epsilon).
\]
The restriction of \(\widetilde X\) over \(V\) is affine, say
\[
    \widetilde V=\operatorname{Spec}\widetilde B.
\]
Since \(\widetilde X\) is flat over \(\Delta\) and has central fiber \(X\), the \(R\)-algebra
\(\widetilde B\) fits into a square-zero extension
\[
    0\longrightarrow B
    \xrightarrow{\iota}
    \widetilde B
    \xrightarrow{q}
    B
    \longrightarrow 0.
\]
Here \(q\) is reduction modulo \(\epsilon\), and \(\iota\) identifies \(B\) with
\(\epsilon\widetilde B\). Explicitly, if \(b\in B\) and \(\widetilde b\in \widetilde B\) is any lift of \(b\),
then
\[
    \iota(b)=\epsilon\widetilde b.
\]
This is independent of the chosen lift because \(\epsilon^2=0\). Thus
\[
    \ker(q)=\epsilon\widetilde B\simeq B,
    \qquad
    (\epsilon\widetilde B)^2=0.
\]

The morphism \(\widetilde\varphi\) gives a ring homomorphism
\[
    \widetilde\varphi^\#:
    A\oplus A\epsilon
    \longrightarrow
    \widetilde B.
\]
Modulo \(\epsilon\), this recovers the original morphism
\[
    \varphi^\#:A\to B.
\]
Since \(\varphi\) factors through \(Y\), the kernel of \(\varphi^\#:A\to B\) is \(I\), and the image
of \(\varphi^\#\) is the summand \(A/I\subset B=A/I\oplus M\).

We first construct the map
\[
    \nu_{\widetilde\varphi}: I/I^2\to B.
\]
If \(a\in I\), then
\[
    q\bigl(\widetilde\varphi^\#(a)\bigr)
    =
    \varphi^\#(a)
    =
    0.
\]
Therefore
\[
    \widetilde\varphi^\#(a)\in \ker(q)=\epsilon\widetilde B\simeq B.
\]
Define \(\lambda(a)\in B\) by the condition
\[
    \iota(\lambda(a))=\widetilde\varphi^\#(a).
\]
If \(a,b\in I\), then
\[
    \widetilde\varphi^\#(a),\widetilde\varphi^\#(b)\in \epsilon\widetilde B.
\]
Hence
\[
    \widetilde\varphi^\#(ab)
    =
    \widetilde\varphi^\#(a)\widetilde\varphi^\#(b)
    \in
    (\epsilon\widetilde B)^2
    =
    0.
\]
Thus \(\lambda(ab)=0\), so \(\lambda\) kills \(I^2\). Therefore \(\lambda\) descends to a homomorphism
\[
    \overline{\lambda}: I/I^2\to B.
\]
This is the local form of \(\nu_{\widetilde\varphi}\).

Now decompose
\[
    B=A/I\oplus M.
\]
Write
\[
    \overline{\lambda}=\nu_1\oplus \nu_2,
\]
where
\[
    \nu_1:I/I^2\to A/I,
    \qquad
    \nu_2:I/I^2\to M.
\]

Let \(\widetilde J\subset A\oplus A\epsilon\) be the ideal of the scheme-theoretic image
\(\widetilde Y\subset W\times\Delta\). Thus
\[
    \widetilde J=\ker(\widetilde\varphi^\#).
\]
We now compute \(\widetilde J\). Let \(a+a'\epsilon\in A\oplus A\epsilon\). If
\(a+a'\epsilon\in\widetilde J\), then reducing modulo \(\epsilon\) gives
\[
    \varphi^\#(a)=0,
\]
so \(a\in I\).

Assume now that \(a\in I\). Then
\[
    \widetilde\varphi^\#(a)=\iota(\lambda(a)).
\]
Also,
\[
    \widetilde\varphi^\#(a'\epsilon)
    =
    \epsilon\widetilde\varphi^\#(a')
    =
    \iota(\varphi^\#(a')).
\]
Therefore
\[
    \widetilde\varphi^\#(a+a'\epsilon)
    =
    \iota\bigl(\lambda(a)+\varphi^\#(a')\bigr).
\]
Since \(\iota\) is injective, this vanishes if and only if
\[
    \lambda(a)+\varphi^\#(a')=0
    \quad\text{in }B.
\]
Equivalently, writing \(B=A/I\oplus M\), this condition becomes
\[
    \nu_1(\overline a)+\overline{a'}=0
    \quad\text{in }A/I,
\]
and
\[
    \nu_2(\overline a)=0
    \quad\text{in }M,
\]
where \(\overline a\) denotes the class of \(a\) in \(I/I^2\), and
\(\overline{a'}\) denotes the class of \(a'\) in \(A/I\).

Thus
\[
    \widetilde J
    =
    \left\{
        a+a'\epsilon
        \ \middle|\
        a\in I,\ 
        \nu_1(\overline a)+\overline{a'}=0,\ 
        \nu_2(\overline a)=0
    \right\}.
\]

The ideal \(J\subset A\) of the central fiber \(\widetilde Y_0\subset W\) is
\[
    J=
    \frac{\widetilde J+(\epsilon)}{(\epsilon)}
    \subset A.
\]
Equivalently, \(a\in J\) if and only if there exists \(a'\in A\) such that
\[
    a+a'\epsilon\in \widetilde J.
\]
By the formula for \(\widetilde J\), this is equivalent to
\[
    a\in I,\qquad \nu_2(\overline a)=0,
\]
together with the existence of \(a'\in A\) satisfying
\[
    \overline{a'}=-\nu_1(\overline a).
\]
The last condition is automatic, because every element of \(A/I\) has a lift to \(A\).
Therefore
\[
    J
    =
    \{a\in I\mid \nu_2(\overline a)=0\}.
\]
Equivalently,
\[
    J=
    \ker\left(
        I\to I/I^2\xrightarrow{\nu_2}M
    \right).
\]
Sheafifying this local computation gives
\[
    I_{\widetilde Y_0/Z}
    =
    \ker\left(
        I_{Y/Z}
        \longrightarrow
        I_{Y/Z}/I_{Y/Z}^2
        \xrightarrow{\nu_2}
        \mathcal E
    \right).
\]

Since the map factors through \(I/I^2\), we have \(I^2\subset I_{\widetilde Y_0/Z}\).
Also \(I_{\widetilde Y_0/Z}\subset I\). Hence
\[
    I^2\subset I_{\widetilde Y_0/Z}\subset I,
\]
which is equivalent to
\[
    Y\subset \widetilde Y_0\subset Y^{(1)}_Z.
\]
This proves the theorem.
\end{proof}

\section{The groups and sheaves related to deformations}

\begin{lemma}\label{lem:T1-relative-differentials}
Let $\pi:X\to Y$ be a morphism locally of finite type between noetherian schemes. Assume that $X$ is $S_1$ and $\pi$ is generically smooth at every generic point of $X$.  
Then the natural map
$
\mathcal E xt_X^1(\Omega_{X/Y},\mathcal O_X)
\longrightarrow
\mathcal T^1(X/Y)
$
is an isomorphism, where
$
\mathcal T^1(X/Y):=\mathcal E xt_X^1(\mathbb L_{X/Y},\mathcal O_X).
$
In particular, the conclusion holds if $\pi$ is finite and generically étale.
\end{lemma}

\begin{proof}
The question is local on $X$. Choose a local factorization
$
X\hookrightarrow P\to Y
$
with $P\to Y$ smooth, and let $I\subset \mathcal O_P$ be the ideal sheaf of $X$ in $P$. The cotangent complex $\mathbb L_{X/Y}$ is represented by the two-term complex
$
\left[
I/I^2
\xrightarrow{d}
\Omega_{P/Y}\otimes_{\mathcal O_P}\mathcal O_X
\right],
$
where $I/I^2$ is placed in degree $-1$ and $\Omega_{P/Y}\otimes_{\mathcal O_P}\mathcal O_X$ is placed in degree $0$.

Set $F:=I/I^2,$ $G:=\Omega_{P/Y}\otimes_{\mathcal O_P}\mathcal O_X,$ $
N:=\operatorname{Im}(F\to G),$ and $K:=\ker(F\to G).$
Then there is an exact sequence
$$
0\to N\to G\to \Omega_{X/Y}\to0.
$$
Since $G$ is locally free, applying $\mathcal H om_X(-,\mathcal O_X)$ gives
$$
\mathcal E xt_X^1(\Omega_{X/Y},\mathcal O_X) =
\operatorname{coker}
\left(
\mathcal H om_X(G,\mathcal O_X)
\to
\mathcal H om_X(N,\mathcal O_X)
\right).
$$

On the other hand, by the above two-term presentation of $\mathbb L_{X/Y}$, we have
$$
\mathcal T^1(X/Y) =
\operatorname{coker}
\left(
\mathcal H om_X(G,\mathcal O_X)
\to
\mathcal H om_X(F,\mathcal O_X)
\right).
$$


Moreover, from the exact sequence
$
0\to K\to F\to N\to0
$
we get
$$
0\to
\mathcal H om_X(N,\mathcal O_X)
\to
\mathcal H om_X(F,\mathcal O_X)
\to
\mathcal H om_X(K,\mathcal O_X).
$$
Suppose $X$ is $S_1$ and $\pi$ is generically smooth at every generic point of $X$. Then the conormal map
$
I/I^2\to \Omega_{P/Y}\otimes_{\mathcal O_P}\mathcal O_X
$
is injective at every generic point of $X$. Hence $K$ is supported away from the generic points of $X$. Since $X$ is $S_1$, the sheaf $\mathcal O_X$ has no nonzero subsheaves supported away from the generic points. Therefore
$
\mathcal H om_X(K,\mathcal O_X)=0.
$
Therefore
$
\mathcal H om_X(N,\mathcal O_X)
\simeq
\mathcal H om_X(F,\mathcal O_X).
$
Substituting this into the two cokernel descriptions above gives
$
\mathcal E xt_X^1(\Omega_{X/Y},\mathcal O_X)
\simeq
\mathcal T^1(X/Y).
$

In particular, if $\pi$ is finite and generically étale, then $f$ is generically smooth, so the same argument applies.
\end{proof}

\begin{lemma}\label{lem:omega-double-cover-trace-zero}
Let $k$ be a field with $\operatorname{char}k\neq 2$, and let $\pi:X\to Y$ be a finite morphism of noetherian $k$-schemes. Assume that, as an $\mathcal{O}_Y$-algebra, one has a decomposition
$$
\pi_{*}\mathcal{O}_X=\mathcal{O}_Y\oplus\mathcal{E},
$$
where the product of two local sections of $\mathcal{E}$ lands in $\mathcal{O}_Y$. Equivalently, the multiplication on the trace-zero summand is given by an $\mathcal{O}_Y$-bilinear map
$$
\mu:\mathcal{E}\otimes_{\mathcal{O}_Y}\mathcal{E}\longrightarrow\mathcal{O}_Y.
$$
Let $J\subset\mathcal{O}_X$ be the ideal generated by the image of the natural map
$$
\pi^*\mathcal{E}\longrightarrow\mathcal{O}_X,
$$
and set $R:=V(J)$.

Then there is a natural exact sequence
$$
0\longrightarrow\mathcal{K}\longrightarrow\Omega_{X/Y}
\longrightarrow
\mathcal{O}_R\otimes_{\mathcal{O}_X}\pi^*\mathcal{E}
\longrightarrow 0.
$$
Moreover, over the open subset of $Y$ where $\mathcal{E}$ is locally free of rank $1$, the map
$$
\Omega_{X/Y}
\longrightarrow
\mathcal{O}_R\otimes_{\mathcal{O}_X}\pi^*\mathcal{E}
$$
is an isomorphism. Consequently,
$$
\operatorname{Supp}\mathcal{K}
\subset
\pi^{-1}\bigl(\operatorname{NLF}(\mathcal{E})\bigr),
$$
where $\operatorname{NLF}(\mathcal{E})$ denotes the non-locally-free locus of $\mathcal{E}$.

Assume further that $\mathcal{K}$ has finite length and that
$$
\operatorname{depth}\mathcal{O}_{X,x}\geq 2
$$
for every $x\in\operatorname{Supp}\mathcal{K}$. Then the natural map induced by the above exact sequence gives an isomorphism
$$
\mathcal{E}xt_X^1
\left(
\mathcal{O}_R\otimes_{\mathcal{O}_X}\pi^*\mathcal{E},
\mathcal{O}_X
\right)
\simeq
\mathcal{E}xt_X^1(\Omega_{X/Y},\mathcal{O}_X).
$$
In particular, this conclusion holds if $\mathcal{K}$ has finite length and $X$ is $S_2$ along $\operatorname{Supp}\mathcal{K}$, with all points of $\operatorname{Supp}\mathcal{K}$ having codimension at least $2$ in $X$.
\end{lemma}

\begin{proof}
The construction is local on $Y$. Let $U=\operatorname{Spec}A\subset Y$ be an affine open subset, and write $\pi^{-1}(U)=\operatorname{Spec}C$. Put
$$
M:=\mathcal{E}(U).
$$
Then the algebra $\pi_*\mathcal{O}_X$ over $U$ is
$$
C=A\oplus M
$$
as an $A$-module. The multiplication on $C$ is determined by the $A$-bilinear map
$$
\mu:M\otimes_A M\longrightarrow A.
$$
Thus, for $m,n\in M$, the product of $m$ and $n$ inside $C$ is
$$
m\cdot n=\mu(m,n)\in A.
$$

Let
$$
B:=\operatorname{Sym}*A(M).
$$
The inclusion $M\hookrightarrow C$ induces an $A$-algebra homomorphism
$$
\varphi:B\longrightarrow C.
$$
This map is surjective, because $C=A\oplus M$ is generated as an $A$-algebra by $M$. Let
$$
I:=\ker(\varphi).
$$
The conormal sequence for the closed embedding $\operatorname{Spec}C\hookrightarrow\operatorname{Spec}B$ over $\operatorname{Spec}A$ gives
$$
I/I^2
\longrightarrow
\Omega_{B/A}\otimes_B C
\longrightarrow
\Omega_{C/A}
\longrightarrow 0.
$$
Since $B=\operatorname{Sym}_A(M)$, one has the standard identification
$$
\Omega*{B/A}\simeq B\otimes_A M.
$$
After tensoring with $C$, this becomes
$$
\Omega_{B/A}\otimes_B C\simeq C\otimes_A M.
$$
Thus the conormal map may be viewed as a map
$$
\delta:I/I^2\longrightarrow C\otimes_A M.
$$

We first describe the ideal $I$. For every $m,n\in M$, the element
$$
mn-\mu(m,n)\in B
$$
maps to zero in $C$, because the product of $m$ and $n$ in $C$ is $\mu(m,n)$. Hence $mn-\mu(m,n)\in I$. Let $I_2\subset B$ be the ideal generated by all such elements:
$$
I_2:=\bigl(mn-\mu(m,n)\mid m,n\in M\bigr).
$$
We claim that $I=I_2$.

Set
$$
Q:=B/I_2.
$$
The map $\varphi:B\to C$ factors through $Q$, giving an $A$-algebra map
$$
Q\longrightarrow C.
$$
Conversely, define an $A$-linear map
$$
\psi:C=A\oplus M\longrightarrow Q
$$
by
$$
\psi(a+m)=\overline{a}+\overline{m}.
$$
This is well-defined because $C=A\oplus M$ is a direct sum as an $A$-module. We check that $\psi$ is multiplicative. Let $a,b\in A$ and $m,n\in M$. Then
$$
\psi(a+m)\psi(b+n)
(\overline{a}+\overline{m})(\overline{b}+\overline{n}).
$$
Expanding in $Q$, we obtain
$$
\psi(a+m)\psi(b+n)
\overline{ab}+\overline{an}+\overline{bm}+\overline{mn}.
$$
By the defining relation in $Q$, one has
$$
\overline{mn}=\overline{\mu(m,n)}.
$$
Therefore
$$
\psi(a+m)\psi(b+n)
\overline{ab}+\overline{an}+\overline{bm}+\overline{\mu(m,n)}.
$$
On the other hand, in $C$ one has
$$
(a+m)(b+n)=ab+an+bm+\mu(m,n).
$$
Applying $\psi$ gives
$$
\psi\bigl((a+m)(b+n)\bigr)
\overline{ab}+\overline{an}+\overline{bm}+\overline{\mu(m,n)}.
$$
Hence
$$
\psi(a+m)\psi(b+n)=\psi\bigl((a+m)(b+n)\bigr).
$$
Thus $\psi$ is an $A$-algebra homomorphism. The maps $Q\to C$ and $C\to Q$ are inverse to each other, because both compositions are the identity on $A$ and on $M$, and $Q$ is generated as an $A$-algebra by the image of $M$. Therefore $Q\simeq C$, and hence $I=I_2$.

Now compute the conormal map on the generators of $I$. Since $\mu(m,n)\in A$, its relative differential over $A$ is zero. Hence
$$
d\bigl(mn-\mu(m,n)\bigr)=m,dn+n,dm.
$$
Under the identification
$$
\Omega_{B/A}\otimes_B C\simeq C\otimes_A M,
$$
this corresponds to
$$
m\otimes n+n\otimes m.
$$
Let $J_C\subset C$ be the ideal generated by the image of $M\subset C$. Since $m,n\in M\subset J_C$, we have
$$
m,dn+n,dm\in J_C\cdot(C\otimes_A M).
$$
Since the elements $mn-\mu(m,n)$ generate $I$, it follows that
$$
\operatorname{Im}(\delta)\subset J_C\cdot(C\otimes_A M).
$$
Therefore the quotient map
$$
C\otimes_A M\longrightarrow (C/J_C)\otimes_A M
$$
kills the image of $\delta$. Hence it descends to a natural surjection
$$
\Omega_{C/A}\twoheadrightarrow (C/J_C)\otimes_A M.
$$

Now $(C/J_C)\otimes_A M$ is naturally identified with the restriction of $\mathcal{O}_R\otimes_{\mathcal{O}_X}\pi^*\mathcal{E}$ to $\pi^{-1}(U)$. Indeed,
$$
(C/J_C)\otimes_C(C\otimes_A M)\simeq (C/J_C)\otimes_A M.
$$
Thus the local maps glue and give a natural surjection
$$
\Omega_{X/Y}
\twoheadrightarrow
\mathcal{O}_R\otimes_{\mathcal{O}_X}\pi^*\mathcal{E}.
$$
Let $\mathcal{K}$ be its kernel. This gives the exact sequence
$$
0\longrightarrow\mathcal{K}\longrightarrow\Omega_{X/Y}
\longrightarrow
\mathcal{O}_R\otimes_{\mathcal{O}_X}\pi^*\mathcal{E}
\longrightarrow 0.
$$

It remains to check the assertion on the locus where $\mathcal{E}$ is locally free of rank $1$. On this locus, after shrinking $U$ if necessary, we may write $M=Az$ for a generator $z$. Then
$$
B=\operatorname{Sym}_A(M)\simeq A[T],
$$
where $T$ corresponds to $z$. If $f:=\mu(z,z)\in A$, then
$$
C\simeq A[T]/(T^2-f).
$$
Identifying the image of $T$ in $C$ with $z$, the ideal $J_C$ is $(z)$. The conormal map is given by
$$
T^2-f\longmapsto d(T^2-f)=2z,dz.
$$
Since $\operatorname{char}k\neq 2$, the image of the conormal map is exactly $(z)\cdot C,dz$. Therefore
$$
\Omega_{C/A}\simeq C,dz/(z\cdot C,dz)\simeq (C/(z))\otimes_A Az.
$$
This is precisely $(C/J_C)\otimes_A M$. Hence the surjection
$$
\Omega_{C/A}\twoheadrightarrow (C/J_C)\otimes_A M
$$
is an isomorphism wherever $\mathcal{E}$ is locally free of rank $1$. Therefore
$$
\operatorname{Supp}\mathcal{K}
\subset
\pi^{-1}\bigl(\operatorname{NLF}(\mathcal{E})\bigr).
$$

Now assume that $\mathcal{K}$ has finite length and that $\operatorname{depth}\mathcal{O}_{X,x}\geq 2$ for every $x\in\operatorname{Supp}\mathcal{K}$. Then
$$
\mathcal{H}om_X(\mathcal{K},\mathcal{O}_X)=0
$$
and
$$
\mathcal{E}xt_X^1(\mathcal{K},\mathcal{O}_X)=0.
$$
Indeed, locally this follows from the standard depth vanishing statement: if $A'$ is a noetherian local ring, $N$ is a finite-length $A'$-module, and $\operatorname{depth}A'\geq 2$, then
$$
\operatorname{Ext}_{A'}^i(N,A')=0
$$
for $i=0,1$.

Applying $\mathcal{H}om_X(-,\mathcal{O}-X)$ to
$$
0\longrightarrow\mathcal{K}\longrightarrow\Omega_{X/Y}
\longrightarrow
\mathcal{O}_R\otimes_{\mathcal{O}_X}\pi^*\mathcal{E}
\longrightarrow 0
$$
and using the two vanishings above, the long exact sequence of sheaf Exts gives
$$
\mathcal{E}xt_X^1
\left(
\mathcal{O}_R\otimes_{\mathcal{O}_X}\pi^*\mathcal{E},
\mathcal{O}_X
\right)
\simeq
\mathcal{E}xt_X^1(\Omega_{X/Y},\mathcal{O}_X).
$$
Finally, if $\mathcal{K}$ has finite length and $X$ is $S_2$ along $\operatorname{Supp}\mathcal{K}$, with all points of $\operatorname{Supp}\mathcal{K}$ having codimension at least $2$ in $X$, then $\operatorname{depth}\mathcal{O}_{X,x}\geq 2$ for every $x\in\operatorname{Supp}\mathcal{K}$, so the preceding argument applies.
\end{proof}

\begin{lemma}\label{lem:differentials-symmetric-algebra}(\textcolor{blue}{Explanatory})
Let $A$ be a ring, let $M$ be an $A$-module, and set
$$
B:=\operatorname{Sym}_A(M).
$$
Then there is a natural isomorphism of $B$-modules
$$
\Omega_{B/A}\simeq B\otimes_A M.
$$
Under this isomorphism, the universal derivation
$$
d:B\longrightarrow \Omega_{B/A}
$$
sends an element $m\in M\subset B$ to
$$
dm=1\otimes m.
$$

\end{lemma}

\begin{proof}
Choose a presentation of $M$ by a free $A$-module:
$$
N\xrightarrow{\alpha} F\longrightarrow M\longrightarrow 0,
$$
where $F$ is free. Let
$$
P:=\operatorname{Sym}_A(F).
$$
Since $F$ is free, $P$ is a polynomial algebra over $A$, possibly in infinitely many variables. Therefore
$$
\Omega_{P/A}\simeq P\otimes_A F.
$$
The surjection $F\to M$ induces a surjection of $A$-algebras
$$
P=\operatorname{Sym}_A(F)\longrightarrow \operatorname{Sym}_A(M)=B.
$$
Let $J$ be its kernel. Then $J$ is the ideal of $P$ generated by the image of $N$ in $F\subset P$. Equivalently,
$$
J=(\alpha(n)\mid n\in N).
$$
Thus
$$
B=P/J.
$$

Apply the conormal sequence to the quotient $P\to B$ over $A$. We get an exact sequence
$$
J/J^2
\longrightarrow
\Omega_{P/A}\otimes_P B
\longrightarrow
\Omega_{B/A}
\longrightarrow 0.
$$
Using $\Omega_{P/A}\simeq P\otimes_A F$, this becomes
$$
J/J^2
\longrightarrow
B\otimes_A F
\longrightarrow
\Omega_{B/A}
\longrightarrow 0.
$$

We now identify the image of the first map. For $n\in N$, the element $\alpha(n)\in F\subset P$ lies in $J$. Its differential is
$$
d(\alpha(n))=1\otimes \alpha(n)\in P\otimes_A F.
$$
After tensoring with $B$, this becomes
$$
1\otimes \alpha(n)\in B\otimes_A F.
$$
More generally, every element of $J$ is a finite sum of terms $p_i\alpha(n_i)$, with $p_i\in P$ and $n_i\in N$. Its differential is a sum of terms
$$
d\bigl(p_i\alpha(n_i)\bigr) =
p_id(\alpha(n_i))+\alpha(n_i)dp_i.
$$
After tensoring with $B=P/J$, the second term vanishes because $\alpha(n_i)\in J$. Therefore the image of $d\bigl(p_i\alpha(n_i)\bigr)$ in $B\otimes_A F$ is
$$
\overline{p_i}\otimes \alpha(n_i).
$$
It follows that the image of
$$
J/J^2\longrightarrow B\otimes_A F
$$
is exactly the image of the natural map
$$
B\otimes_A N\longrightarrow B\otimes_A F.
$$

Therefore the conormal sequence gives
$$
\Omega_{B/A}
\simeq
\operatorname{coker}\left(B\otimes_A N\longrightarrow B\otimes_A F\right).
$$
On the other hand, tensoring the right-exact sequence
$$
N\longrightarrow F\longrightarrow M\longrightarrow 0
$$
with $B$ over $A$ gives a right-exact sequence
$$
B\otimes_A N
\longrightarrow
B\otimes_A F
\longrightarrow
B\otimes_A M
\longrightarrow 0.
$$
Hence
$$
\operatorname{coker}\left(B\otimes_A N\longrightarrow B\otimes_A F\right)
\simeq
B\otimes_A M.
$$
Combining the preceding identifications, we obtain
$$
\Omega_{B/A}\simeq B\otimes_A M.
$$

Finally, under this isomorphism, if $m\in M\subset B$, choose a lift $f\in F$ of $m$. In $P$, one has
$$
df=1\otimes f.
$$
Passing to $B$ gives
$$
dm=1\otimes m.
$$
This proves the claim.
\end{proof}


\section{Deformations of double covers of cones over a rational normal curve}

\subsection{Canonical double covers of rational normal cones}

Let $e>4$, and let $Y=S(0,e)\subset \mathbb{P}^{e+1}$ be the rational normal cone of degree $e$. Equivalently, $Y\simeq \mathbb{P}(1,1,e)$. We denote by $p\in Y$ the vertex. The point $p$ is the unique singular point of $Y$, and analytically it is a cyclic quotient singularity of type $\frac{1}{e}(1,1)$.

Let $F$ denote the generator of the class group of $Y$, so that $\mathcal{O}_{Y}(F)=\mathcal{O}_{\mathbb{P}(1,1,e)}(1)$. The hyperplane class of the embedding $Y\subset \mathbb{P}^{e+1}$ is $H=eF$, or equivalently
$$
\mathcal{O}_{Y}(H)\simeq \mathcal{O}_{\mathbb{P}(1,1,e)}(e).
$$
The canonical sheaf of $Y$ is
$$
\omega_{Y}\simeq \mathcal{O}_{\mathbb{P}(1,1,e)}(-(e+2)).
$$

We consider a finite double cover $\pi:X\to Y$ whose trace decomposition is
$$
\pi_{*}\mathcal{O}_{X}=\mathcal{O}_{Y}\oplus \mathcal{E}.
$$
For the cover to be canonical with respect to the hyperplane class $H$, we require $\omega_{X}\simeq \pi^*\mathcal{O}_{Y}(H)$. The double-cover canonical formula gives
$$
\omega_{X}\simeq \pi^*(\omega_{Y}\otimes \mathcal{E}^{\vee}),
$$
where $\mathcal{E}^{\vee}:=\mathcal{H}om_{Y}(\mathcal{E},\mathcal{O}_{Y})$. Therefore the condition $\omega_{X}\simeq \pi^*\mathcal{O}_{Y}(H)$ forces $\omega_{Y}\otimes \mathcal{E}^{\vee}\simeq \mathcal{O}_{Y}(H)$, and hence
$$
\mathcal{E}\simeq \omega_{Y}(-H).
$$
Substituting the expressions above, we obtain
$$
\mathcal{E}
\simeq
\mathcal{O}_{\mathbb{P}(1,1,e)}(-(e+2)-e)
= \mathcal{O}_{\mathbb{P}(1,1,e)}(-(2e+2)).
$$

Thus the double cover is determined by the algebra $\mathcal{O}_{Y}\oplus\mathcal{E}$, together with a multiplication map
$$
\mathcal{E}\otimes_{\mathcal{O}_{Y}}\mathcal{E}\longrightarrow \mathcal{O}_{Y}.
$$
which is the same as a multiplication map 
$$
(\mathcal{E}\otimes_{\mathcal{O}_{Y}}\mathcal{E})^{\vee \vee}\longrightarrow \mathcal{O}_{Y}.
$$
Equivalently, this multiplication is determined by a section of the reflexive square $(\mathcal{E}^{\vee}\otimes_{\mathcal{O}_{Y}}\mathcal{E}^{\vee})^{\vee\vee}$. Since $\mathcal{E}^{\vee}\simeq \mathcal{O}_{\mathbb{P}(1,1,e)}(2e+2)$, we have
$$
(\mathcal{E}^{\vee}\otimes_{\mathcal{O}_{Y}}\mathcal{E}^{\vee})^{\vee\vee}
\simeq
\mathcal{O}_{\mathbb{P}(1,1,e)}(4e+4).
$$
Hence the branch data are given by a section
$$
s\in H^0\left(Y,\mathcal{O}_{\mathbb{P}(1,1,e)}(4e+4)\right).
$$
We denote the corresponding branch divisor by $B:=V(s)$. Thus
$$
B\in \left|\mathcal{O}_{\mathbb{P}(1,1,e)}(4e+4)\right|.
$$

Since $e>4$, every weighted homogeneous polynomial of degree $4e+4$ in coordinates of weights $1,1,e$ vanishes at the vertex $p=[0:0:1]$. Indeed, a monomial involving only the weight-$e$ coordinate would have degree divisible by $e$, while $4e+4\not\equiv 0\pmod e$ for $e>4$. Therefore every monomial of weighted degree $4e+4$ contains at least one of the weight-$1$ coordinates, and hence vanishes at $p$. Consequently, $p\in B$.

The trace-zero sheaf $\mathcal{E}=\mathcal{O}_{\mathbb{P}(1,1,e)}(-(2e+2))$ is rank-one reflexive. It is locally free at the vertex precisely when $e$ divides $2e+2$, equivalently when $e$ divides $2$. Since $e>4$, this does not occur. Thus $\mathcal{E}$ is not locally free at $p$, and the double cover is non-flat over the vertex in the sense of Alexeev--Pardini.

Away from $p$, the sheaf $\mathcal{E}$ is locally free, and the cover is locally given by an equation $z^2=f$. At the vertex, however, the cover must be understood through the reflexive algebra $\mathcal{O}_{Y}\oplus\mathcal{E}$ rather than through a line-bundle total space.

With this setup, the canonical divisor of $X$ satisfies $K_{X}=\pi^*H$. Since $H^2=e$ on the rational normal cone, one obtains
$$
K_{X}^2=2H^2=2e.
$$
Moreover,
$$
p_g(X)=h^0(X,\omega_{X})=h^0(Y,\mathcal{O}_{Y}(H))=e+2.
$$
Thus
$$
K_{X}^2=2e=2p_g(X)-4.
$$



\begin{proposition}\label{prop:hom-conormal-trace-zero-cone}
Let $Y=S(0,e)\subset\mathbb{P}^{e+1}$ be the rational normal cone of degree $e$, and let $H=\mathcal{O}_Y(1)$ denote the hyperplane class of this embedding. Equivalently, identify
$$
Y\simeq \mathbb{P}(1,1,e),
$$
so that
$$
H\simeq \mathcal{O}_{\mathbb{P}(1,1,e)}(e).
$$
Let $\mathcal{I}\subset\mathcal{O}_{\mathbb{P}^{e+1}}$ be the ideal sheaf of $Y$. Set
$$
\mathcal{E}:=\omega_Y(-H).
$$
Then
$$
\operatorname{Hom}_Y(\mathcal{I}/\mathcal{I}^2,\mathcal{E})=0.
$$
\end{proposition}

\begin{proof}
The homogeneous ideal of the rational normal cone $Y=S(0,e)\subset\mathbb{P}^{e+1}$ is generated by quadrics. Therefore there is a surjection
$$
\mathcal{O}_Y(-2H)^{\oplus N}\twoheadrightarrow \mathcal{I}/\mathcal{I}^2
$$
for some integer $N$. Applying $\operatorname{Hom}_Y(-,\mathcal{E})$ gives an injection
$$
\operatorname{Hom}_Y(\mathcal{I}/\mathcal{I}^2,\mathcal{E})
\hookrightarrow
\operatorname{Hom}_Y(\mathcal{O}_Y(-2H)^{\oplus N},\mathcal{E}).
$$
The right-hand side is
$$
H^0(Y,\mathcal{E}(2H))^{\oplus N}.
$$
Thus it is enough to show that
$$
H^0(Y,\mathcal{E}(2H))=0.
$$

Now
$$
\mathcal{E}(2H)=\omega_Y(-H+2H)=\omega_Y(H).
$$
Under the identification $Y\simeq\mathbb{P}(1,1,e)$, one has
$$
\omega_Y\simeq\mathcal{O}_{\mathbb{P}(1,1,e)}(-(e+2)).
$$
Since $H\simeq\mathcal{O}_{\mathbb{P}(1,1,e)}(e)$, it follows that
$$
\mathcal{E}(2H)
\simeq
\mathcal{O}_{\mathbb{P}(1,1,e)}(-(e+2)+e)
= \mathcal{O}_{\mathbb{P}(1,1,e)}(-2).
$$
But
$$
H^0(\mathbb{P}(1,1,e),\mathcal{O}_{\mathbb{P}(1,1,e)}(-2))=0.
$$
Therefore
$$
H^0(Y,\mathcal{E}(2H))=0.
$$
Hence
$$
\operatorname{Hom}_Y(\mathcal{I}/\mathcal{I}^2,\mathcal{E})=0.
$$
\end{proof}



\section{Computations on a weighted projective space}

\subsection{Theorems proven in this section so far}

\begin{theorem}
Let $Y = \PP(1,1,m,m)$ with $m \geq 4$. Then
\[
\mathcal{T}^1_Y:=\mathcal{E}xt^1(\Omega_Y, \OO_Y) \;\cong\; \OO_{\PP^1}(1)^{\oplus(m-3)},
\]
supported on the singular locus $\Sigma = \{x_0 = x_1 = 0\} \cong \PP^1$.
\end{theorem}

\begin{proposition}
Let
\[
Y=\PP(1,1,m,m),
\qquad
\Sigma=\Sing(Y)\simeq \PP^1.
\]
Then
\[
\mathcal{T}^2_Y:=\mathcal{E}xt^2(\Omega_Y,\OO_Y)
\cong
\OO_{\PP^1}(2)^{\oplus (m-1)(m-3)}.
\]
\end{proposition}


\begin{proposition}
Let
\[
Y=\mathbb{P}(1,1,m,m).
\]
Assume \(m\geq 4\). Then \(Y\) has obstructed deformations.
\end{proposition}



\begin{set-up}\label{embedding_of_WPS_in_PN}
Let $Y=\mathbb P(a_0,\dots,a_r)$ be a weighted projective space over an algebraically closed field,
of dimension $r\ge 2$, so that $Y$ is normal. Fix integers $m>0$ and $n>0$.
Let $Z=\mathbb P^N$ and let $i:Y\hookrightarrow Z$ be a closed embedding induced by a very ample line bundle
$L:=i^*\mathcal O_{\mathbb P^N}(1)\cong \mathcal O_Y(m)$.
Set $\mathcal E:=\bigoplus_i\mathcal O_Y(-n_i)$.
\end{set-up}
\begin{lemma}\label{lemma:hom_euler_wps}
In Set-up \ref{embedding_of_WPS_in_PN},

\begin{enumerate}
\item[\textup{(1)}] There is a short exact sequence of $\mathcal O_Y$--modules
\begin{equation}\label{eq:hom_euler_exact}
0 \longrightarrow \mathcal O_Y(-n)\longrightarrow
\mathcal O_Y(m-n)^{\oplus (N+1)}
\longrightarrow
\mathcal H\!om\!\bigl(i^*\Omega_{\mathbb P^N},\,\mathcal O_Y(-n)\bigr)
\longrightarrow 0.
\end{equation}

\item[\textup{(2)}] Taking global sections yields an exact sequence
\begin{equation}\label{eq:hom_euler_globalsections}
0 \longrightarrow H^0\!\bigl(Y,\mathcal O_Y(-n)\bigr)\longrightarrow
H^0\!\bigl(Y,\mathcal O_Y(m-n)\bigr)^{\oplus (N+1)}
\longrightarrow
\operatorname{Hom}\bigl(i^*\Omega_{\mathbb P^N},\mathcal O_Y(-n)\bigr)
\longrightarrow 0.
\end{equation}
In particular,
\[
\operatorname{Hom}\bigl(i^*\Omega_{\mathbb P^N},\mathcal O_Y(-n)\bigr)
\cong
H^0\!\bigl(Y,\mathcal O_Y(m-n)\bigr)^{\oplus (N+1)}.
\]

\item[\textup{(3)}] If $m<n$, then $m-n<0$, hence $H^0\!\bigl(Y,\mathcal O_Y(m-n)\bigr)=0$ and therefore
\[
\operatorname{Hom}\bigl(i^*\Omega_{\mathbb P^N},\mathcal O_Y(-n)\bigr)=0.
\]
\end{enumerate}

\end{lemma}

\begin{proof}
Restrict the Euler sequence on $\mathbb P^N$ to $Y$:
\[
0\to i^*\Omega_{\mathbb P^N}\to \mathcal O_Y(-m)^{\oplus (N+1)}\to \mathcal O_Y\to 0.
\]
Apply $\mathcal H\!om(-,\mathcal O_Y(-n))$.
Since $\mathcal O_Y$ and $\mathcal O_Y(-m)$ are locally free, $\mathcal E\!xt^1(\mathcal O_Y,\mathcal O_Y(-n))=0$,
and we obtain \eqref{eq:hom_euler_exact}. Taking $H^0$ and using $r\ge 2$ gives
$H^1(Y,\mathcal O_Y(-n))=0$, so the long exact sequence truncates to \eqref{eq:hom_euler_globalsections},
yielding (2). Statement (3) follows immediately from $m-n<0$.
\end{proof}

\begin{lemma}\label{lem:Ext1_pullback_Omega_vanish_WPS}
Let $Y=\mathbb P(a_0,\dots,a_r)$ be a weighted projective space over an algebraically closed field,
with $\dim Y=r\ge 3$. Let $i:Y\hookrightarrow \mathbb P^N$ be a closed embedding such that
$i^*\mathcal O_{\mathbb P^N}(1)\cong \mathcal O_Y(m)$ for some $m>0$.
Let $\mathcal E$ be a finite direct sum of twists
\[
\mathcal E \;=\;\bigoplus_{i=1}^s \mathcal O_Y(-n_i).
\]
Then
\[
\operatorname{Ext}^1\!\bigl(i^*\Omega_{\mathbb P^N},\,\mathcal E\bigr)\;=\;0.
\]
Equivalently, for every integer $-n_i$ one has
\[
\operatorname{Ext}^1\!\bigl(i^*\Omega_{\mathbb P^N},\,\mathcal O_Y(-n_i)\bigr)\;=\;0.
\]
\end{lemma}

\begin{proof}
It suffices to prove the statement for $\mathcal E=\mathcal O_Y(k)$, since $\operatorname{Ext}^1$ commutes with finite
direct sums in the second argument. Because $i^*\Omega_{\mathbb P^N}$ is locally free on $Y$, we have
\[
\operatorname{Ext}^1\!\bigl(i^*\Omega_{\mathbb P^N},\,\mathcal O_Y(k)\bigr)
\;\cong\;
H^1\!\bigl(Y,\mathcal H\!om(i^*\Omega_{\mathbb P^N},\mathcal O_Y(k))\bigr)
\;\cong\;
H^1\!\bigl(Y,\,i^*T_{\mathbb P^N}\otimes \mathcal O_Y(k)\bigr).
\]
Restrict the dual Euler sequence on $\mathbb P^N$ to $Y$ and twist by $\mathcal O_Y(k)$:
\[
0\longrightarrow \mathcal O_Y(k)
\longrightarrow \mathcal O_Y(k+m)^{\oplus (N+1)}
\longrightarrow i^*T_{\mathbb P^N}\otimes \mathcal O_Y(k)
\longrightarrow 0.
\]
Taking cohomology yields an exact segment
\[
H^1\!\bigl(Y,\mathcal O_Y(k+m)\bigr)^{\oplus (N+1)}
\longrightarrow
H^1\!\bigl(Y,i^*T_{\mathbb P^N}\otimes \mathcal O_Y(k)\bigr)
\longrightarrow
H^2\!\bigl(Y,\mathcal O_Y(k)\bigr).
\]
For weighted projective space $Y$ of dimension $r\ge 3$, one has
\[
H^i\!\bigl(Y,\mathcal O_Y(t)\bigr)=0\qquad\text{for all integers $t$ and all }0<i<r,
\]
see for example \cite[Theorem~3.2.4(iii)]{DolgachevWPS}.
Hence $H^1(Y,\mathcal O_Y(k+m))=0$ and $H^2(Y,\mathcal O_Y(k))=0$, so the middle term vanishes:
\[
H^1\!\bigl(Y,i^*T_{\mathbb P^N}\otimes \mathcal O_Y(k)\bigr)=0.
\]
Therefore $\operatorname{Ext}^1(i^*\Omega_{\mathbb P^N},\mathcal O_Y(k))=0$, as claimed.
\end{proof}

\begin{proposition}\label{embedded_ribbons_abstract_ribbons}

Under the conditions of Set-up \ref{embedding_of_WPS_in_PN}, if $m < n_i$ for all $i$, then we have 

\begin{equation*}
  \operatorname{Hom}(\mathcal{I}/\mathcal{I}^2, \mathcal{E}) \cong \operatorname{Ext}^1(\Omega_Y, \mathcal{E})  
\end{equation*}

\begin{proof}
    This follows from \Cref{lem:Ext1_pullback_Omega_vanish_WPS}, \Cref{eq:hom_euler_globalsections} and \Cref{counting_reflexive_ribbons}.
\end{proof}

    
\end{proposition}

\begin{theorem}\label{thm:main}
Let $Y = \PP(1,1,m,m)$ with $m \geq 4$. Then
\[
\mathcal{E}xt^1(\Omega_Y, \OO_Y) \;\cong\; \OO_{\PP^1}(1)^{\oplus(m-3)},
\]
supported on the singular locus $\Sigma = \{x_0 = x_1 = 0\} \cong \PP^1$.
\end{theorem}
 
\subsection{Notation}
 
Let $Y = \Proj\, k[x_0, x_1, x_2, x_3]$ where $\deg x_0 = \deg x_1 = 1$ and $\deg x_2 = \deg x_3 = m$, with $m \geq 4$. The singular locus of $Y$ is
\[
\Sigma = \{x_0 = x_1 = 0\} \cong \PP(m,m) \cong \PP^1.
\]
Every point $p = [0:0:x_2:x_3] \in \Sigma$ has stabilizer $\ZZ/m$ under the group action defining $Y$. The generator $\zeta$ (a primitive $m$-th root of unity) acts on the transverse coordinates as $(x_0, x_1) \mapsto (\zeta x_0, \zeta x_1)$, so the transverse singularity type is $\frac{1}{m}(1,1)$.
 
Let $V$ denote the affine cone over the rational normal curve $\nu_m(\PP^1) \subset \PP^m$. Its coordinate ring is
\[
R = k[z_0, \ldots, z_m] \,\big/\, I_2\!\begin{pmatrix} z_0 & z_1 & \cdots & z_{m-1} \\ z_1 & z_2 & \cdots & z_m \end{pmatrix},
\]
with the standard grading $\deg z_i = 1$, so that $V = \Spec\, R$. The ideal of $2 \times 2$ minors is denoted $I \subset S := k[z_0, \ldots, z_m]$, with generators $f_{ij} = z_i z_{j+1} - z_{i+1} z_j$ of degree $2$. The singularity $\frac{1}{m}(1,1)$ is isomorphic to the germ of $V$ at the vertex.
 
\subsection{The charts}
 
\begin{lemma}\label{lem:charts}
The singular locus $\Sigma$ is covered by two affine open sets $\Spec(B_t)$ and $\Spec(B_s)$, where
\[
B_t = R \otimes_k k[t], \qquad B_s = R \otimes_k k[s],
\]
so that $\Spec(B_t) \cong V \times \AAA^1_t$ and $\Spec(B_s) \cong V \times \AAA^1_s$. The transition isomorphism $B_t[t^{-1}] \xrightarrow{\;\sim\;} B_s[s^{-1}]$ is given by
\begin{equation}\label{eq:transition_rings}
b_k = \frac{a_k}{t}, \qquad s = \frac{1}{t},
\end{equation}
where $a_k$ and $b_k$ denote the images of the generators $z_k$ in $B_t$ and $B_s$ respectively.
\end{lemma}
 
\begin{proof}
\textbf{Chart $\Spec(B_t)$ (where $x_2 \neq 0$):} Define
\[
a_k = \frac{x_0^{m-k}\, x_1^k}{x_2}, \quad 0 \leq k \leq m, \qquad t = \frac{x_3}{x_2}.
\]
All expressions have degree $0$ in the weighted grading (numerator and denominator both have degree $m$ for $a_k$; likewise $t$ has degree $0$). The coordinate ring is
\[
B_t = k[a_0, \ldots, a_m, t] \,\big/\, (a_i a_{j+1} - a_{i+1} a_j)_{i < j} \;\cong\; R \otimes_k k[t],
\]
since $t$ is algebraically independent of the $a_k$.
 
\medskip
 
\textbf{Chart $\Spec(B_s)$ (where $x_3 \neq 0$):} Define
\[
b_k = \frac{x_0^{m-k}\, x_1^k}{x_3}, \quad 0 \leq k \leq m, \qquad s = \frac{x_2}{x_3}.
\]
Similarly, $B_s = k[b_0, \ldots, b_m, s] / (b_i b_{j+1} - b_{i+1} b_j) \cong R \otimes_k k[s]$.
 
\medskip
 
\textbf{Transition on the overlap:} Since
\[
b_k = \frac{x_0^{m-k}x_1^k}{x_3} = \frac{x_0^{m-k}x_1^k}{x_2} \cdot \frac{x_2}{x_3} = \frac{a_k}{t}, \qquad s = \frac{x_2}{x_3} = \frac{1}{t},
\]
the transition is given by~\eqref{eq:transition_rings}. While each chart is individually a product, the gluing mixes the cone coordinates with the base parameter.
\end{proof}
 
\subsection{$T^1$ of a product with a smooth factor}
 
\begin{lemma}\label{lem:product}
Let $A$ be a $k$-algebra and $C$ the coordinate ring of a smooth affine $k$-scheme. Then
\[
T^1_{A \otimes C} \;\cong\; T^1_A \otimes_k C.
\]
\end{lemma}
 
\begin{proof}
Choose a presentation $A = S/I$ where $S = k[z_0, \ldots, z_m]$. Then $A \otimes C = (S \otimes C)/(I \otimes C)$, where $S \otimes C$ is a polynomial ring over $C$ (smooth, since $\Spec(C)$ is smooth).
 
Now $T^1_A$ is computed from the exact sequence
\[
\Der_k(S, A) \xrightarrow{\;\alpha\;} \Hom_A(I/I^2, A) \to T^1_A \to 0.
\]
Since $C$ is flat over $k$, tensoring with $C$ preserves exactness:
\[
\Der_k(S, A) \otimes_k C \xrightarrow{\;\alpha \otimes 1\;} \Hom_A(I/I^2, A) \otimes_k C \to T^1_A \otimes_k C \to 0.
\]
On the other hand, the analogous sequence for $A \otimes C$ is
\[
\Der_k(S \otimes C, A \otimes C) \xrightarrow{\;\alpha'\;} \Hom_{A \otimes C}\!\big((I \otimes C)/(I \otimes C)^2,\, A \otimes C\big) \to T^1_{A \otimes C} \to 0.
\]
We have $\Hom_{A \otimes C}((I \otimes C)/(I \otimes C)^2, A \otimes C) \cong \Hom_A(I/I^2, A) \otimes_k C$ by flatness. For the derivation module, $\Der_k(S \otimes C, A \otimes C)$ decomposes as $\Der_k(S, A) \otimes_k C \;\oplus\; \Hom_k(\Omega_C, A \otimes C)$, but the second summand maps to zero under $\alpha'$ since derivations of $C$ annihilate the generators of $I$. Therefore $T^1_{A \otimes C} = T^1_A \otimes_k C$.
\end{proof}
 
\begin{corollary}\label{cor:product}
$T^1_{B_t} = T^1_R \otimes_k k[t]$ and $T^1_{B_s} = T^1_R \otimes_k k[s]$.
\end{corollary}

\subsection{$T^i$ of a product with a smooth factor}
\begin{lemma}\label{lemma:cotangent-cohomology-polynomial-extension}
Let $k$ be a field, let $B$ be a $k$-algebra, and set
\[
C=B[t]\cong B\otimes_k k[t].
\]
Then for every $i\geq 1$ there is a natural isomorphism
\[
T^i_{C/k}\cong T^i_{B/k}\otimes_B C.
\]
Equivalently,
\[
\Ext^i_C(L_{C/k},C)\cong \Ext^i_B(L_{B/k},B)\otimes_B C
\qquad \text{for all } i\geq 1.
\]
\end{lemma}

\begin{proof}
Set
\[
R=k,\qquad A=B,\qquad S=k[t],\qquad C=A\otimes_R S.
\]
Since $S=k[t]$ is flat over $R=k$, the $R$-algebras $A$ and $S$ are Tor-independent. Therefore, by the tensor-product formula for the cotangent complex, there is an isomorphism in $D(C)$
\[
L_{C/R}\cong
\bigl(L_{A/R}\otimes_A^{\mathbf L} C\bigr)
\oplus
\bigl(L_{S/R}\otimes_S^{\mathbf L} C\bigr)
\]
(see \cite[\href{https://stacks.math.columbia.edu/tag/09DA}{Tag 09DA}]{StacksProject}).

Now $S=k[t]$ is smooth over $k$, so
\[
L_{S/R}=L_{k[t]/k}\simeq \Omega_{k[t]/k}[0]
\]
by \cite[\href{https://stacks.math.columbia.edu/tag/08R5}{Tag 08R5}]{StacksProject}.
Since $\Omega_{k[t]/k}$ is a free $k[t]$-module of rank one, it follows that
\[
L_{S/R}\otimes_S^{\mathbf L} C \simeq \Omega_{k[t]/k}\otimes_{k[t]} C
\]
is a free $C$-module placed in degree $0$. Hence
\[
\Ext^i_C\!\left(L_{S/R}\otimes_S^{\mathbf L} C,\; C\right)=0
\qquad \text{for all } i\geq 1.
\]
Therefore, for $i\geq 1$,
\[
\Ext^i_C(L_{C/R},C)\cong
\Ext^i_C\!\left(L_{A/R}\otimes_A^{\mathbf L} C,\; C\right).
\]

Choose a free $A$-resolution $P^\bullet\to L_{A/R}$. Since $C$ is flat over $A$, the complex
\[
P^\bullet\otimes_A C
\]
is a free $C$-resolution of $L_{A/R}\otimes_A^{\mathbf L} C$. Thus
\[
\Ext^i_C\!\left(L_{A/R}\otimes_A^{\mathbf L} C,\; C\right)
=
H^i\!\left(\Hom_C(P^\bullet\otimes_A C,\; C)\right).
\]
Using adjunction,
\[
\Hom_C(P^\bullet\otimes_A C,\; C)\cong \Hom_A(P^\bullet,C).
\]
Since each $P^n$ is a free $A$-module and $C$ is flat over $A$, we obtain
\[
\Hom_A(P^\bullet,C)\cong \Hom_A(P^\bullet,A)\otimes_A C.
\]
Taking cohomology and using again that $C$ is flat over $A$, we get
\[
H^i\!\left(\Hom_A(P^\bullet,A)\otimes_A C\right)
\cong
H^i\!\left(\Hom_A(P^\bullet,A)\right)\otimes_A C.
\]
But
\[
H^i\!\left(\Hom_A(P^\bullet,A)\right)
=
\Ext^i_A(L_{A/R},A)
=
T^i_{A/R}.
\]
Hence, for all $i\geq 1$,
\[
T^i_{C/R}
=
\Ext^i_C(L_{C/R},C)
\cong
T^i_{A/R}\otimes_A C.
\]
Substituting $A=B$, $R=k$, and $C=B[t]$ gives
\[
T^i_{C/k}\cong T^i_{B/k}\otimes_B C
\qquad \text{for all } i\geq 1,
\]
as claimed.
\end{proof}
 
\subsection{Pinkham's computation of $T^1_R$}
 
\begin{lemma}\label{lem:pinkham}
Let $R$ be the coordinate ring of the cone over the rational normal curve $X = \nu_m(\PP^1) \subset \PP^m$, with $m \geq 4$. Then $T^1_R$, as a graded $R$-module, is concentrated in degree $-1$:
\[
(T^1_R)_d = \begin{cases} k^{m-3} & \text{if } d = -1, \\ 0 & \text{otherwise.} \end{cases}
\]
\end{lemma}
 
\begin{proof}
The rational normal curve $X = \nu_m(\PP^1) \subset \PP^m$ has $\OO_{\PP^m}(1)|_X = \OO_{\PP^1}(m)$, tangent bundle $\Theta_X = \OO_{\PP^1}(2)$, and normal bundle $N_{X/\PP^m} = \OO_{\PP^1}(m+2)^{\oplus(m-1)}$.
 
From the exact sequence $0 \to \Theta_X \to \Theta_{\PP^m}|_X \to N_{X/\PP^m} \to 0$, twisted by $\OO_X(d) = \OO_{\PP^1}(md)$, we obtain the long exact sequence in cohomology:
\begin{equation}\label{eq:les}
H^0(\Theta_{\PP^m}|_X(d)) \xrightarrow{\;\delta_d\;} H^0(N_X(d)) \to H^1(\Theta_X(d)) \to H^1(\Theta_{\PP^m}|_X(d)).
\end{equation}
By Pinkham's theorem, for $d < 0$:
\[
(T^1_R)_d = \coker(\delta_d).
\]
 
\medskip
 
\textbf{Case $d = -1$.} Twist the Euler sequence $0 \to \OO_X \to \OO_X(1)^{m+1} \to \Theta_{\PP^m}|_X \to 0$ by $\OO_X(-1) = \OO_{\PP^1}(-m)$:
\[
0 \to \OO_{\PP^1}(-m) \to \OO_{\PP^1}^{\;m+1} \to \Theta_{\PP^m}|_X(-1) \to 0.
\]
Since $H^1(\OO_{\PP^1}^{\;m+1}) = 0$, the long exact sequence gives $H^1(\Theta_{\PP^m}|_X(-1)) = 0$. From~\eqref{eq:les}, the cokernel of $\delta_{-1}$ is therefore
\[
(T^1_R)_{-1} \;=\; H^1(\Theta_X(-1)) \;=\; H^1(\OO_{\PP^1}(2 - m)) \;=\; m - 3 \quad\text{for } m \geq 4.
\]
 
\medskip
 
\textbf{Case $d \leq -2$.} We have $N_X(d) = \OO_{\PP^1}(m + 2 + md)^{\oplus(m-1)}$, and $m + 2 + md \leq m + 2 - 2m = 2 - m < 0$ for $m \geq 3$. So $H^0(N_X(d)) = 0$, hence $(T^1_R)_d = 0$.
 
\medskip
 
\textbf{Case $d \geq 0$.} For $d \geq 0$, $\Theta_X(d) = \OO_{\PP^1}(2 + md)$ has $H^1 = 0$, so the long exact sequence~\eqref{eq:les} shows $\delta_d$ is surjective. Furthermore, by the Euler sequence twisted by $\OO_X(d)$, $H^1(\Theta_{\PP^m}|_X(d)) = 0$, and $H^1(N_X(d)) = H^1(\OO_{\PP^1}(m+2+md))^{\oplus(m-1)} = 0$. Therefore $(T^1_R)_d = 0$.
\end{proof}
 
\subsection{The grading on $T^1_{B_t}$}
 
\begin{lemma}\label{lem:grading}
The module $T^1_{B_t} = T^1_R \otimes_k k[t]$ carries a natural grading from the cone variables $a_0, \ldots, a_m$ (with $\deg a_i = 1$ and $\deg t = 0$). In this grading, $T^1_{B_t}$ is a free $k[t]$-module of rank $m - 3$ concentrated in cone-degree $-1$.
 
Concretely, fix a $k$-basis $\phi_1, \ldots, \phi_{m-3}$ of $(T^1_R)_{-1}$. Each $\phi_\alpha$ is a class in the quotient $\Hom_R(I/I^2, R)_{-1} / \im(\Der_k(S,R)_{-1})$, represented by an $R$-module map $I/I^2 \to R$ sending each generator $f_{ij}$ to a linear form $\ell^\alpha_{ij}(z_0, \ldots, z_m) \in R_1$, subject to the syzygy constraints and modulo the image of constant derivations.
 
A general element of $T^1_{B_t}$ is
\[
\sum_{\alpha=1}^{m-3} c_\alpha(t)\, \phi_\alpha, \qquad c_\alpha(t) \in k[t],
\]
which sends $f_{ij}(a)$ to the linear form $\sum_\alpha c_\alpha(t)\, \ell^\alpha_{ij}(a_0, \ldots, a_m) \in R_1 \otimes k[t]$.
 
Similarly, a general element of $T^1_{B_s}$ is $\sum_\alpha d_\alpha(s)\, \phi_\alpha$ with $d_\alpha(s) \in k[s]$, sending $g_{ij}(b)$ to $\sum_\alpha d_\alpha(s)\, \ell^\alpha_{ij}(b_0, \ldots, b_m)$.
\end{lemma}
 
\begin{proof}
Since $B_t = R \otimes_k k[t]$ with $R$ graded by $\deg a_i = 1$ and $t$ in degree $0$, and $T^1_{B_t} = T^1_R \otimes_k k[t]$ by Corollary~\ref{cor:product}, the module $T^1_{B_t}$ inherits the $R$-grading on the first factor. By Lemma~\ref{lem:pinkham}, $T^1_R$ is concentrated in cone-degree~$-1$, so:
\[
(T^1_{B_t})_d = (T^1_R)_d \otimes_k k[t] = \begin{cases} k^{m-3} \otimes_k k[t] \cong k[t]^{m-3} & \text{if } d = -1, \\ 0 & \text{otherwise.} \end{cases} \qedhere
\]
\end{proof}
 
\subsection{Proof of Theorem~\ref{thm:main}}
 
By Lemmas~\ref{lem:charts}--\ref{lem:grading}, the sheaf $\Tt^1_Y = \mathcal{E}xt^1(\Omega_Y, \OO_Y)$ is supported on $\Sigma \cong \PP^1$ and is obtained by gluing the free $k[t]$-module $T^1_{B_t} \cong k[t]^{m-3}$ and the free $k[s]$-module $T^1_{B_s} \cong k[s]^{m-3}$ over the overlap.
 
With the notation of Lemma~\ref{lem:grading}, a section of $\Tt^1_Y$ over $\Spec(B_t)$ has the form $\sum_\alpha c_\alpha(t)\, \phi_\alpha$, which sends
\[
f_{ij}(a) \;\longmapsto\; \sum_\alpha c_\alpha(t)\, \ell^\alpha_{ij}(a_0, \ldots, a_m).
\]
 
On the overlap, we re-express this in $\Spec(B_s)$-coordinates using the transition $a_k = t\, b_k$ and $s = 1/t$ from Lemma~\ref{lem:charts}. Write $f_{ij}(a) = a_i a_{j+1} - a_{i+1}a_j$ and $g_{ij}(b) = b_i b_{j+1} - b_{i+1}b_j$ for the generators of $I$ in the two charts. Then:
\begin{itemize}
\item The generators transform as $f_{ij}(a) = t^2\, g_{ij}(b)$, since $f_{ij}$ is quadratic in the cone variables.
\item Each linear form transforms as $\ell^\alpha_{ij}(a_0, \ldots, a_m) = t \cdot \ell^\alpha_{ij}(b_0, \ldots, b_m)$, since $\ell^\alpha_{ij}$ is homogeneous of degree $1$ in the cone variables.
\end{itemize}
The same $R$-module map, viewed as acting on the $\Spec(B_s)$-generators $g_{ij}(b)$, gives:
\[
g_{ij}(b) = \frac{f_{ij}(a)}{t^2} \;\longmapsto\; \frac{1}{t^2} \sum_\alpha c_\alpha(t)\, \ell^\alpha_{ij}(a) = \frac{1}{t^2} \sum_\alpha c_\alpha(t) \cdot t \cdot \ell^\alpha_{ij}(b) = \frac{1}{t}\sum_\alpha c_\alpha(t)\, \ell^\alpha_{ij}(b).
\]
Comparing with the $\Spec(B_s)$-expression $g_{ij}(b) \mapsto \sum_\alpha d_\alpha(s)\, \ell^\alpha_{ij}(b)$, we find the transition
\begin{equation}\label{eq:transition}
d_\alpha(s) \;=\; \frac{c_\alpha(t)}{t} \;=\; s \cdot c_\alpha(1/s), \qquad \alpha = 1, \ldots, m-3.
\end{equation}
 
The transition~\eqref{eq:transition} is diagonal in the basis $\phi_1, \ldots, \phi_{m-3}$: there is no mixing between different $\alpha$'s. Each component transforms by $c_\alpha(t) \mapsto s \cdot c_\alpha(1/s)$, which is the standard transition function for $\OO_{\PP^1}(1)$. (Recall: a section of $\OO_{\PP^1}(1)$ is given by $c(t) \in k[t]$ on the chart $\{s \neq 0\}$ and $d(s) \in k[s]$ on the chart $\{t \neq 0\}$, with $d(s) = s \cdot c(1/s)$ on the overlap.) Therefore:
\[
\mathcal{E}xt^1(\Omega_Y, \OO_Y) \;\cong\; \OO_{\PP^1}(1)^{\oplus(m-3)}. \qed
\]
 
\begin{remark}
The transition~\eqref{eq:transition} is diagonal because all basis elements $\phi_\alpha$ live in the same graded piece $(T^1_R)_{-1}$, and the substitution $a_k = t\, b_k$ acts on $R_1$ by a uniform scalar factor $t$. The syzygy constraints and the derivation quotient defining $(T^1_R)_{-1}$ are homogeneous, hence $\mathbb{G}_m$-equivariant, so they do not introduce any mixing between the basis vectors.
\end{remark}
 
\subsection{Consequences and special cases}
 
\begin{itemize}
\item $h^0(\Tt^1_Y) = 2(m-3)$: the dimension of the space of first-order deformations smoothing the singularities of~$Y$.
\item $h^1(\Tt^1_Y) = 0$: no obstructions at the $H^1(\Tt^1)$ level in the local-to-global spectral sequence $E_2^{p,q} = H^p(\Tt^q_Y) \Rightarrow T^{p+q}(Y)$.
\end{itemize}
 
\begin{remark}[Special cases]\leavevmode
\begin{itemize}
\item For $m = 3$, $(T^1_R)_{-1} = H^1(\OO_{\PP^1}(-1)) = 0$, so $\Tt^1_Y = 0$ and $\PP(1,1,3,3)$ has no local-to-global smoothing deformations.
\item For $m = 2$, the singularity $\frac{1}{2}(1,1) \cong A_1$ is a hypersurface $\{z_0 z_2 - z_1^2 = 0\} \subset \AAA^3$. Here $T^1_R$ is $1$-dimensional, concentrated in cone-degree $-2$ (a constant perturbation $z_0 z_2 - z_1^2 = c$). The transition gives $d(s) = s^2 \cdot c(1/s)$, yielding $\Tt^1_Y = \OO_{\PP^1}(2)$ with $h^0 = 3$.
\end{itemize}
\end{remark}

\subsection{Calculating $T_Y^2$}

\begin{proposition}\label{prop:TY2-wps-1122m}
Let
\[
Y=\PP(1,1,m,m),
\qquad
\Sigma=\Sing(Y)\simeq \PP^1.
\]
Then
\[
\mathcal{T}^2_Y:=\mathcal{E}xt^2(\Omega_Y,\OO_Y)
\cong
\OO_{\PP^1}(2)^{\oplus (m-1)(m-3)}.
\]
\end{proposition}

\begin{proof}
Let
\[
Y=\Proj k[x_0,x_1,y_0,y_1],
\qquad
\deg x_0=\deg x_1=1,\quad \deg y_0=\deg y_1=m.
\]
Its singular locus is
\[
\Sigma=(x_0=x_1=0)\simeq \PP(m,m)\simeq \PP^1.
\]

Consider the two standard affine charts
\[
U_0=(y_0\neq 0),\qquad U_1=(y_1\neq 0).
\]
On \(U_0\), set
\[
t=\frac{y_1}{y_0},
\qquad
a_i=\frac{x_0^{m-i}x_1^i}{y_0}
\quad (0\leq i\leq m).
\]
Then
\[
U_0\simeq \Spec(B[t]),
\]
where
\[
B=
k[a_0,\dots,a_m]\Big/
I_2\!\begin{pmatrix}
a_0&a_1&\cdots&a_{m-1}\\
a_1&a_2&\cdots&a_m
\end{pmatrix}
\]
is the affine cone over the rational normal curve of degree \(m\).

Similarly, on \(U_1\), set
\[
s=\frac{y_0}{y_1},
\qquad
b_i=\frac{x_0^{m-i}x_1^i}{y_1}
\quad (0\leq i\leq m),
\]
so that
\[
U_1\simeq \Spec(B[s]).
\]
where by a slight abuse of notation,
\[
B=
k[b_0,\dots,b_m]\Big/
I_2\!\begin{pmatrix}
b_0&b_1&\cdots&b_{m-1}\\
b_1&b_2&\cdots&b_m
\end{pmatrix}
\]
On the overlap \(U_0\cap U_1\), one has
\[
s=t^{-1},
\qquad
b_i=t^{-1}a_i
\quad (0\leq i\leq m).
\]

By Lemma~\ref{lemma:cotangent-cohomology-polynomial-extension}, for every \(i\geq 1\),
\[
T^i_{B[t]/k}\cong T^i_{B/k}\otimes_B B[t],
\qquad
T^i_{B[s]/k}\cong T^i_{B/k}\otimes_B B[s].
\]
In particular,
\[
T^2(B[t])\cong T^2(B)\otimes_B B[t],
\qquad
T^2(B[s])\cong T^2(B)\otimes_B B[s].
\]

Now apply Altmann--Stevens to the cone \(B\) over the rational normal curve.
By \cite[\S4.2(2)]{AltmannStevens}, the only nonzero multigraded pieces of \(T^2(B)\) are
those of multidegree \(R=[k,2]\), and their dimensions are
\[
\dim T^2(-R)=
\begin{cases}
k-2,& R=[k,2],\ 2\leq k\leq m-1,\\[2mm]
m-3,& R=[m,2],\\[2mm]
2m-k-2,& R=[k,2],\ m+1\leq k\leq 2m-2,
\end{cases}
\]
with zero otherwise. Moreover, by \cite[Proposition~4.3(2)]{AltmannStevens},
\[
T^n(-R)=0 \qquad \text{unless } \operatorname{ht}(R)=n.
\]
Hence \(T^2(B)\) is concentrated in ordinary graded degree \(-2\). Therefore, as a graded
\(B\)-module,
\[
T^2(B)\cong k(-2)^{\oplus r},
\]
where
\[
r=\dim_k T^2(B)
=
\sum_{k=2}^{m-1}(k-2)+(m-3)+\sum_{k=m+1}^{2m-2}(2m-k-2).
\]
A direct computation gives
\[
r=\frac{(m-3)(m-2)}{2}+(m-3)+\frac{(m-3)(m-2)}{2}=(m-1)(m-3).
\]
Thus
\[
T^2(B)\cong k(-2)^{\oplus (m-1)(m-3)}.
\]

Since the irrelevant ideal of \(B\) annihilates \(T^2(B)\), we obtain
\[
T^2(B[t])
\cong
k[t]\otimes_k T^2(B),
\qquad
T^2(B[s])
\cong
k[s]\otimes_k T^2(B).
\]
Therefore
\[
\mathcal{T}^2_Y|_{U_0}\cong \OO_{\Sigma\cap U_0}^{\oplus (m-1)(m-3)},
\qquad
\mathcal{T}^2_Y|_{U_1}\cong \OO_{\Sigma\cap U_1}^{\oplus (m-1)(m-3)}.
\]

It remains to determine the transition function on \(U_0\cap U_1\).
Because \(b_i=t^{-1}a_i\), a homogeneous class of ordinary degree \(-2\) transforms by the
factor \(t^{-2}\). Since all of \(T^2(B)\) lies in degree \(-2\), the transition matrix on
\(U_0\cap U_1\) is \(t^{-2}I\). This is precisely the transition function for
\(\OO_{\PP^1}(2)\). Hence
\[
\mathcal{T}^2_Y
\cong
\OO_{\PP^1}(2)^{\oplus (m-1)(m-3)}.
\]

This proves the proposition.
\end{proof}

\begin{proposition}\label{prop:wps-1122m-obstructed}
Let
\[
Y=\mathbb{P}(1,1,m,m).
\]
Assume \(m\geq 4\). Then \(Y\) has obstructed deformations.
\end{proposition}

\begin{proof}
Let
\[
\Sigma=\operatorname{Sing}(Y)\simeq \mathbb{P}^{1}.
\]
On the standard affine charts
\[
U_{0}=(y_{0}\neq 0),\qquad U_{1}=(y_{1}\neq 0),
\]
one has
\[
U_{0}\simeq \operatorname{Spec}(B[t]),\qquad U_{1}\simeq \operatorname{Spec}(B[s]),
\]
where
\[
B=
k[a_{0},\dots,a_{m}]
\Big/
I_{2}\!\begin{pmatrix}
a_{0}&a_{1}&\cdots&a_{m-1}\\
a_{1}&a_{2}&\cdots&a_{m}
\end{pmatrix}
\]
is the affine cone over the rational normal curve of degree \(m\). On the overlap
\(U_{0}\cap U_{1}\) the coordinates satisfy
\[
s=t^{-1},\qquad b_{i}=t^{-1}a_{i}\qquad (0\leq i\leq m).
\]

It is classical that the cone over the rational normal curve of degree \(m\) is obstructed
for \(m\geq 4\); see Pinkham \cite[Chapter~8]{Pinkham}. A convenient explicit reference is
\cite{CilibertoLopezMiranda}. Choose a first-order deformation class
\[
\zeta_{B}\in T^{1}_{B/k}
\]
which does not lift to second order. Let
\[
A_{1}=k[\varepsilon]/(\varepsilon^{2}),\qquad A_{2}=k[\varepsilon]/(\varepsilon^{3}),
\]
and let \(B_{1}\) be a flat \(A_{1}\)-deformation of \(B\) representing \(\zeta_{B}\).

The induced deformations \(B_{1}[t]\) and \(B_{1}[s]\) determine first-order deformations
of \(U_{0}\) and \(U_{1}\). By the previously established identification
\[
\mathcal{T}^{1}_{Y}\cong \mathcal{O}_{\mathbb{P}^{1}}(1)\otimes_{k}T^{1}_{B/k},
\]
the transition on \(U_{0}\cap U_{1}\) is multiplication by \(t^{-1}=s\), because
\(T^{1}_{B/k}\) is concentrated in degree \(-1\). Therefore the local sections on
\(U_{0}\cap \Sigma\) and \(U_{1}\cap \Sigma\) corresponding to the same vector
\(\zeta_{B}\in T^{1}_{B/k}\) in the two local trivializations satisfy the correct transition
rule on the overlap. Hence they glue to a global section
\[
\sigma\in H^{0}(\mathcal{T}^{1}_{Y}).
\]

By the local-to-global spectral sequence for
\(\operatorname{Ext}^{\bullet}(\Omega_{Y},\mathcal{O}_{Y})\), there is an exact sequence
\[
0\longrightarrow H^{1}(T_{Y})
\longrightarrow \operatorname{Ext}^{1}(\Omega_{Y},\mathcal{O}_{Y})
\longrightarrow H^{0}(\mathcal{T}^{1}_{Y})
\longrightarrow H^{2}(T_{Y}).
\]
Since \(H^{2}(T_{Y})=0\) by the vanishing proved earlier, the middle map is surjective.
Choose
\[
\eta\in \operatorname{Ext}^{1}(\Omega_{Y},\mathcal{O}_{Y})
\]
mapping to \(\sigma\).

We claim that \(\eta\) is obstructed. Suppose, on the contrary, that \(\eta\) lifts to a
second-order deformation of \(Y\). Restricting this second-order deformation to the affine chart
\(U_{0}\), we obtain a flat \(A_{2}\)-algebra \(C_{2}\) together with a surjection
\[
q\colon C_{2}\twoheadrightarrow B_{1}[t]
\]
whose kernel is square-zero. Since \(B_{1}[t]\) is naturally an \(A_{2}[t]\)-algebra via the
quotient map \(A_{2}[t]\twoheadrightarrow A_{1}[t]\), and since the polynomial algebra
\(A_{2}[t]\) is formally smooth over \(A_{2}\), the map \(A_{2}[t]\to B_{1}[t]\) lifts across
\(q\) to an \(A_{2}\)-algebra map
\[
\widetilde{\phi}\colon A_{2}[t]\longrightarrow C_{2}
\]
(see \cite[Tag~00TH]{StacksProject}). Let
\[
\widetilde{t}:=\widetilde{\phi}(t)\in C_{2}.
\]

Set
\[
Q:=C_{2}/(\widetilde{t}).
\]
We will show that \(Q\) is a flat \(A_{2}\)-deformation of \(B\) whose first-order truncation
is \(B_{1}\), which contradicts the choice of \(\zeta_{B}\).

We first prove that multiplication by \(\widetilde{t}\) on \(C_{2}\) is injective. Let
\(x\in C_{2}\) satisfy \(\widetilde{t}x=0\). Reducing modulo \(\varepsilon\), we obtain
\[
t\overline{x}=0 \qquad \text{in } B[t].
\]
Since multiplication by \(t\) is injective on the polynomial ring \(B[t]\), it follows that
\(\overline{x}=0\), hence \(x=\varepsilon x_{1}\) for some \(x_{1}\in C_{2}\). Then
\[
0=\widetilde{t}x=\varepsilon\,\widetilde{t}x_{1}.
\]
Because \(C_{2}\) is flat over \(A_{2}\), one has
\[
\varepsilon m=0 \Longrightarrow m\in \varepsilon^{2}C_{2},
\qquad
\varepsilon^{2}m=0 \Longrightarrow m\in \varepsilon C_{2}.
\]
Therefore \(\widetilde{t}x_{1}\in \varepsilon^{2}C_{2}\), and reducing modulo \(\varepsilon\)
again gives \(t\overline{x_{1}}=0\) in \(B[t]\). Hence \(x_{1}\in \varepsilon C_{2}\), so
\(x=\varepsilon^{2}x_{2}\) for some \(x_{2}\in C_{2}\). Repeating the same argument once more,
we obtain \(x\in \varepsilon^{3}C_{2}=0\). Thus multiplication by \(\widetilde{t}\) is injective,
and we have an exact sequence
\[
0\longrightarrow C_{2}\xrightarrow{\cdot\,\widetilde{t}} C_{2}\longrightarrow Q\longrightarrow 0.
\]

Tensoring this exact sequence with \(A_{1}=A_{2}/(\varepsilon^{2})\), and using the flatness of
\(C_{2}\) over \(A_{2}\), we obtain an exact sequence
\[
0\longrightarrow \operatorname{Tor}^{A_{2}}_{1}(A_{1},Q)
\longrightarrow B_{1}[t]\xrightarrow{\cdot\, t} B_{1}[t]
\longrightarrow Q/\varepsilon^{2}Q
\longrightarrow 0.
\]
Since multiplication by \(t\) is injective on the polynomial ring \(B_{1}[t]\), we conclude that
\[
\operatorname{Tor}^{A_{2}}_{1}(A_{1},Q)=0
\qquad\text{and}\qquad
Q/\varepsilon^{2}Q\cong B_{1}[t]/(t)\cong B_{1}.
\]
As \(B_{1}\) is flat over \(A_{1}\), it follows from the local criterion for flatness
\cite[Tag~051K]{StacksProject} that \(Q\) is flat over \(A_{2}\). Finally,
\[
Q/\varepsilon Q \cong B[t]/(t)\cong B.
\]
Hence \(Q\) is a second-order deformation of \(B\) whose first-order truncation is \(B_{1}\).
This contradicts the choice of \(\zeta_{B}\).

Therefore \(\eta\) does not lift to second order. Hence \(Y\) has obstructed deformations.
\end{proof}




Let $Y = \PP(1,1,m,m)$ with $m \geq 1$, and let $\OO_Y(n)$ be a line bundle with $m \mid n$ and $n > 0$. Let $\pi\colon X \to Y$ be the double cover branched along a divisor $B \in |\OO_Y(2n)|$, defined by a section $s \in H^0(Y, \OO_Y(2n))$. We work over a field of characteristic $\neq 2$.
 
\begin{theorem}\label{thm:T2}
With the notation above,
\[
T^2(X/Y) := \Ext^2(\OO_X;\, \LL^\bullet_{X/Y},\, \OO_X) \;=\; H^1(B,\, \OO_B(2n)).
\]
\end{theorem}
 
We prove this in three steps: first we identify the cotangent complex $\LL^\bullet_{X/Y}$, then we compute its local $\mathcal{E}xt$ sheaves, and finally we apply the local-to-global spectral sequence.
 
\bigskip
 
\subsection{Step 1: The cotangent complex $\LL^\bullet_{X/Y}$}
 
Let $V = \Tot(\OO_Y(n))$ be the total space of the line bundle $\OO_Y(n)$, with projection $p\colon V \to Y$. Let $\tau \in H^0(V,\, p^*\OO_Y(n))$ denote the tautological section. The double cover $X$ is recovered as the zero locus of
\[
h \;=\; \tau^2 - p^*s \;\in\; H^0(V,\, p^*\OO_Y(2n)),
\]
giving a factorization
\[
X \xrightarrow{\;g\;} V \xrightarrow{\;p\;} Y
\]
with $\pi = p \circ g$.
 
\begin{lemma}\label{lem:factorization}
The map $p$ is smooth of relative dimension~$1$, and $g$ is a regular embedding of codimension~$1$.
\end{lemma}
 
\begin{proof}
The map $p$ is the structure morphism of a geometric line bundle, hence smooth of relative dimension~$1$.
 
For $g$: the section $h$ defines $X$ as a Cartier divisor in $V$, with $\OO_V(X) = p^*\OO_Y(2n)$. We must show that $h$ is a non-zero-divisor in $\OO_{V,x}$ at every point $x \in X$.
 
Since $Y$ is a weighted projective space, it is normal, and $V = \Tot(\OO_Y(n))$ is a line bundle over a normal variety, hence itself normal. In particular, $V$ is reduced and satisfies Serre's condition $S_2$. Moreover, $Y = \PP(1,1,m,m)$ has at worst cyclic quotient singularities (transverse $\ZZ/m$ along the line $\ell = \{x_0 = x_1 = 0\}$), so $V$ also has quotient singularities. By the Hochster--Roberts theorem \cite{HR}, $V$ is Cohen--Macaulay.
 
In a Cohen--Macaulay local ring, a single element generating a proper ideal of height~$1$ is a non-zero-divisor. Since $h$ vanishes along the codimension-$1$ subvariety $X \subset V$, the ideal $(h)$ has $\height \geq 1$, and $h$ is therefore a non-zero-divisor in $\OO_{V,x}$ at every $x \in X$. Thus $g$ is a codimension-$1$ regular embedding.
\end{proof}
 
We now compute $\LL^\bullet_{X/Y}$ using the following standard result.
 
\begin{theorem}[Lichtenbaum--Schlessinger; Illusie {\cite[II.1.2.4.2, II.2.2.1]{Illusie}}]\label{thm:naive}
Let $X \xrightarrow{g} P \xrightarrow{\mathrm{smooth}} Y$ be a factorization of a morphism $\pi\colon X \to Y$, with $g$ a regular embedding with ideal sheaf $\cJ \subset \OO_P$. Then the cotangent complex $\LL^\bullet_{X/Y}$ is quasi-isomorphic to the naive cotangent complex
\[
\LL^\bullet_{X/Y} \;\simeq\; \Big[\,\cJ/\cJ^2 \xrightarrow{\;\bar{d}\;} \Omega^1_{P/Y}\big|_X\,\Big]
\]
in degrees $[-1, 0]$, where $\bar{d}(\bar{f}) = df\big|_X$ is the restriction of the relative differential. Moreover, $\cJ/\cJ^2$ is locally free (since $g$ is regular), so this is a perfect complex of amplitude $[-1,0]$, and its quasi-isomorphism type is independent of the chosen factorization.
\end{theorem}
 
We identify the terms of the naive cotangent complex in our situation.
 
\begin{enumerate}
\item \textbf{The conormal sheaf.} The ideal of $X$ in $V$ is $\cI_X = h \cdot \OO_V(-X)$ with $\OO_V(X) = p^*\OO_Y(2n)$, so
\[
\cI_X/\cI_X^2 \;=\; \OO_V(-X)\big|_X \;=\; g^*p^*\OO_Y(-2n) \;=\; \pi^*\OO_Y(-2n).
\]
 
\item \textbf{The relative differentials.} The relative cotangent sheaf of the geometric line bundle $V = \Tot(\OO_Y(n))$ over $Y$ is
\[
\Omega^1_{V/Y} \;=\; p^*\OO_Y(-n),
\]
since the vertical tangent direction is $p^*\OO_Y(n)$ and $\Omega^1_{V/Y}$ is its dual. Restricting to $X$:
\[
g^*\Omega^1_{V/Y} \;=\; \pi^*\OO_Y(-n).
\]
 
\item \textbf{The differential.} The relative differential of the defining equation $h = \tau^2 - p^*s$ is
\[
d_{V/Y}(h) \;=\; d_{V/Y}(\tau^2 - p^*s) \;=\; 2\tau\, d\tau,
\]
since $d_{V/Y}(p^*s) = 0$ (the section $s$ is pulled back from the base $Y$). Under the identifications above, the map $\bar{d}\colon \cI_X/\cI_X^2 \to \Omega^1_{V/Y}\big|_X$ becomes multiplication by $2\tau\big|_X$:
\[
2\tau \;\colon\; \pi^*\OO_Y(-2n) \;\longrightarrow\; \pi^*\OO_Y(-n),
\]
where $\tau\big|_X \in H^0(X, \pi^*\OO_Y(n))$ is the section vanishing along the ramification divisor $R \subset X$.
\end{enumerate}
 
\begin{proposition}\label{prop:cotangent}
The cotangent complex of $\pi\colon X \to Y$ is the two-term complex
\[
\LL^\bullet_{X/Y} \;\simeq\; \Big[\,\pi^*\OO_Y(-2n) \xrightarrow{\;2\tau\;} \pi^*\OO_Y(-n)\,\Big]
\]
in degrees $[-1, 0]$.
\end{proposition}
 
\subsection{Step 2: Local $\mathcal{E}xt$ sheaves}
 
We dualize $\LL^\bullet_{X/Y}$ to compute $R\!\cH\!\mathit{om}_{\OO_X}(\LL^\bullet_{X/Y}, \OO_X)$.
 
\begin{lemma}[Duality for bounded complexes of locally free sheaves]\label{lem:dual}
If $E^\bullet = [E^{-1} \xrightarrow{\delta} E^0]$ is a two-term complex of locally free $\OO_X$-modules in degrees $[-1,0]$, then
\[
R\!\cH\!\mathit{om}_{\OO_X}(E^\bullet, \OO_X) \;\simeq\; \Big[\,(E^0)^\vee \xrightarrow{\;\delta^\vee\;} (E^{-1})^\vee\,\Big]
\]
in degrees $[0, 1]$.
\end{lemma}
 
\begin{proof}
This is the standard computation: applying $\cH\!\mathit{om}(-,\OO_X)$ term-by-term to a bounded complex of locally free sheaves reverses the grading and transposes the differential.
\end{proof}
 
Applying this to $\LL^\bullet_{X/Y}$, the dual complex is
\[
R\!\cH\!\mathit{om}(\LL^\bullet_{X/Y},\, \OO_X) \;\simeq\; \Big[\,\pi^*\OO_Y(n) \xrightarrow{\;2\tau\;} \pi^*\OO_Y(2n)\,\Big]
\]
in degrees $[0, 1]$, where the map is multiplication by $2\tau\big|_X$.
 
\begin{proposition}\label{prop:local_ext}
The local $\mathcal{E}xt$ sheaves (cohomology sheaves of the dual complex) are:
\begin{enumerate}
\item[\textup{(i)}] $\cT^0(X/Y) := \mathcal{E}\mathit{xt}^0(\LL^\bullet_{X/Y}, \OO_X) = 0$.
\item[\textup{(ii)}] $\cT^1(X/Y) := \mathcal{E}\mathit{xt}^1(\LL^\bullet_{X/Y}, \OO_X) \cong \OO_B(2n)$.
\item[\textup{(iii)}] $\cT^q(X/Y) = 0$ for $q \geq 2$.
\end{enumerate}
\end{proposition}
 
\begin{proof}
Part (iii) is immediate since the dual complex is concentrated in degrees $[0,1]$.
 
\medskip
 
\noindent\textup{(i)} We have $\cT^0(X/Y) = \ker(2\tau\colon \pi^*\OO_Y(n) \to \pi^*\OO_Y(2n))$. The sheaves $\pi^*\OO_Y(n)$ and $\pi^*\OO_Y(2n)$ are line bundles on $X$, hence torsion-free. In characteristic $\neq 2$, the section $2\tau\big|_X$ is nonzero (since $\tau\big|_X$ vanishes only along the ramification divisor $R$, which is a proper closed subvariety). Multiplication by a nonzero section is injective on a torsion-free module, hence $\cT^0(X/Y) = 0$.
 
\medskip
 
\noindent\textup{(ii)} We have $\cT^1(X/Y) = \coker(2\tau\colon \pi^*\OO_Y(n) \to \pi^*\OO_Y(2n))$. Since $\mathrm{char} \neq 2$, the factor of~$2$ is a unit, and the cokernel equals
\[
\coker\!\big(\tau\colon \pi^*\OO_Y(n) \to \pi^*\OO_Y(2n)\big).
\]
The section $\tau\big|_X \in H^0(X, \pi^*\OO_Y(n))$ cuts out the ramification divisor $R$ as a Cartier divisor, giving the short exact sequence
\[
0 \;\longrightarrow\; \pi^*\OO_Y(n) \;\xrightarrow{\;\tau\;}\; \pi^*\OO_Y(2n) \;\longrightarrow\; \pi^*\OO_Y(2n)\big|_R \;\longrightarrow\; 0.
\]
Since $\pi\big|_R \colon R \xrightarrow{\;\sim\;} B$ is an isomorphism (the double cover is totally ramifien along $R$), we inentify
\[
\cT^1(X/Y) \;\cong\; \pi^*\OO_Y(2n)\big|_R \;\cong\; \OO_B(2n). \qedhere
\]
\end{proof}
 
\subsection{Step 3: The spectral sequence}
 
\begin{theorem}[Local-to-global spectral sequence; cf.\ Wehler {\cite[1.1]{Wehler}}]\label{thm:ss}
For any object $F^\bullet \in D^b(\OO_X\textup{-mod})$ there is a convergent spectral sequence
\[
E_2^{p,q} \;=\; H^p\!\big(X,\, \cH^q(F^\bullet)\big) \;\Longrightarrow\; \mathbb{H}^{p+q}(X,\, F^\bullet).
\]
Applied to $F^\bullet = R\!\cH\!\mathit{om}(\LL^\bullet_{X/Y}, \OO_X)$, this gives
\[
E_2^{p,q} \;=\; H^p\!\big(X,\, \cT^q(X/Y)\big) \;\Longrightarrow\; T^{p+q}(X/Y).
\]
\end{theorem}
 
\begin{proof}[Proof of Theorem~\ref{thm:T2}]
By Proposition~\ref{prop:local_ext}, the only nonzero local $\mathcal{E}xt$ sheaf is $\cT^1(X/Y) \cong \OO_B(2n)$ in degree $q = 1$. Therefore the $E_2$ page of the spectral sequence has a single nonzero column at $q = 1$:
\[
E_2^{p,q} \;=\;
\begin{cases}
H^p(X,\, \OO_B(2n)) & \text{if } q = 1, \\
0 & \text{if } q \neq 1.
\end{cases}
\]
The spectral sequence degenerates at $E_2$ (all differentials are zero for degree reasons), giving
\[
T^n(X/Y) \;=\; E_2^{n-1,\, 1} \;=\; H^{n-1}\!\big(X,\, \OO_B(2n)\big).
\]
Since $\OO_B(2n)$ is a sheaf supported on $R \cong B \subset X$, we have $H^{n-1}(X, \OO_B(2n)) = H^{n-1}(B, \OO_B(2n))$. Setting $n = 2$:
\[
T^2(X/Y) \;=\; H^1(B,\, \OO_B(2n)). \qedhere
\]
\end{proof}
 
\subsection{Vanishing}
 
\begin{corollary}\label{cor:vanishing}
For $Y = \PP(1,1,m,m)$ with $\dim Y = 3$ and a general branch divisor $B \in |\OO_Y(2n)|$, we have $T^2(X/Y) = 0$.
\end{corollary}
 
\begin{proof}
The branch divisor $B$ is a surface in $Y$. By adjunction,
\[
\omega_B \;=\; \big(\omega_Y \otimes \OO_Y(2n)\big)\big|_B.
\]
Since $\omega_Y = \OO_Y(-2 - 2m)$ for $\PP(1,1,m,m)$, we obtain $\omega_B = \OO_B(2n - 2 - 2m)$. Therefore
\[
H^1(B,\, \OO_B(2n)) \;=\; H^1\!\big(B,\, \omega_B \otimes \OO_B(2 + 2m)\big).
\]
The line bundle $\OO_B(2 + 2m)$ is the restriction of $\OO_Y(2+2m)$ to $B$, which is ample on $B$ since $\OO_Y(1)$ is ample on $\PP(1,1,m,m)$. A general member $B \in |\OO_Y(2n)|$ has at worst quotient singularities (by Bertini's theorem for weighted projective spaces). By Kodaira vanishing for varieties with quotient singularities (cf.\ Steenbrink \cite{Steenbrink}, Saito \cite{Saito}),
\[
H^i(B,\, \omega_B \otimes \mathcal{L}) \;=\; 0 \qquad \text{for } i > 0
\]
whenever $\mathcal{L}$ is ample on a projective variety $B$ with at worst quotient singularities. Applying this with $\mathcal{L} = \OO_B(2+2m)$ and $i = 1$, we conclude $H^1(B, \OO_B(2n)) = 0$, hence $T^2(X/Y) = 0$.
\end{proof}
 
\begin{remark}
For completeness, we record the full tangent cohomology of $\pi$:
\[
T^n(X/Y) \;=\; H^{n-1}(B,\, \OO_B(2n)) \qquad \text{for all } n \geq 0.
\]
In particular, $T^0(X/Y) = 0$ (the double cover has no infinitesimal automorphisms over $Y$, as expected for a $\ZZ/2$-cover), $T^1(X/Y) = H^0(B, \OO_B(2n))$ parametrizes first-order deformations of the branch divisor inside the linear system $|\OO_Y(2n)|$, and $T^2(X/Y) = 0$ says these deformations are unobstructed.
\end{remark}
 

\section{Reflexive Ribbons}
 
Throughout, let $Y$ be a reduced connected scheme of finite type over
an algebraically closed field $k$ of characteristic zero,
and let $\cE$ be a coherent $\cO_Y$-module.
 
\begin{definition}[Reflexive ribbon]\label{def:reflexive-ribbon}
A \emph{reflexive ribbon} over $Y$ with conormal sheaf $\cE$ is a
scheme $\Yt$ such that
\begin{enumerate}
  \item[(i)] $\Yt_{\mathrm{red}} = Y$,
  \item[(ii)] $\cI_{Y,\Yt}^{2} = 0$, and
  \item[(iii)] $\cI_{Y,\Yt} \cong \cE$ as $\cO_Y$-modules,
\end{enumerate}
where $\cE$ is a reflexive sheaf on $Y$.
If $\cE$ is locally free, this recovers the
notion of a \emph{rope} in the sense of \cite{Gon05}; if $\cE$ is
furthermore a line bundle, $\Yt$ is a \emph{ribbon} in the sense of
\cite{BE95}.
\end{definition}
 
\begin{construction}[The algebra structure
{\normalfont (after Bayer--Eisenbud \cite[Theorem~1.2]{BE95})}]
\label{con:algebra}
Let
\[
e \colon \quad 0 \longrightarrow \cE
  \overset{j}{\longrightarrow} \cB
  \overset{p}{\longrightarrow} \Omega_Y \longrightarrow 0
\]
be an extension of $\Omega_Y$ by $\cE$, representing a class
$[e] \in \Ext^{1}_{Y}(\Omega_Y, \cE)$.
Define the sheaf of abelian groups
\[
\cO_{\Yt} \;=\; \cO_Y \times_{\Omega_Y} \cB
\;=\; \bigl\{(a,\, x) \in \cO_Y \oplus \cB
  \;\big|\; da = p(x)\bigr\},
\]
where $d \colon \cO_Y \to \Omega_Y$ is the universal derivation. Equip
$\cO_{\Yt}$ with the multiplication
\[
(a_1, x_1)\cdot (a_2, x_2) = (a_1 a_2,\; a_1 x_2 + a_2 x_1).
\]
This is well-defined: since $d$ is a derivation and $p$ is
$\cO_Y$-linear,
\[
d(a_1 a_2) = a_1\, da_2 + a_2\, da_1
  = a_1\, p(x_2) + a_2\, p(x_1) = p(a_1 x_2 + a_2 x_1).
\]
One verifies that $\cO_{\Yt}$ is a sheaf of commutative $k$-algebras.
Setting $\Yt = (Y, \cO_{\Yt})$, we obtain a reflexive ribbon over $Y$
with conormal sheaf $\cE$.
 
Indeed, the kernel of the natural surjection $\cO_{\Yt} \twoheadrightarrow \cO_Y$, $(a,x)\mapsto a$, is
$\{(0, x) \mid x \in \ker(p)\} \cong \ker(p) = j(\cE) \cong \cE$,
and its square is zero since
$(0,x_1)\cdot(0,x_2) = (0,\, 0)$.
 
%We emphasize that this construction requires no local splitting of
%the extension $e$. The pullback, the multiplication, and the
%verification that $\cI^2 = 0$ are all formal consequences of the
%Leibniz rule and $\cO_Y$-linearity of $p$.
\end{construction}
 
\begin{proposition}[Classification]\label{prop:classification}
Let $Y$ be a reduced connected scheme, proper over $k$, and let
$\cE$ be a reflexive sheaf on $Y$.
\begin{enumerate}
  \item[\textup{(1)}] For every class
    $[e] \in \Ext^{1}_{Y}(\Omega_Y, \cE)$, Construction
    \ref{con:algebra} produces a reflexive ribbon $\Yt$ over $Y$ with
    conormal sheaf $\cE$ and extension class $[e_{\Yt}] = [e]$.
  \item[\textup{(2)}] Conversely, every reflexive ribbon $\Yt$ over
    $Y$ with conormal sheaf $\cE$ determines, via its restricted
    cotangent sequence
    \[
    0 \longrightarrow \cE \longrightarrow \Omega_{\Yt}\big|_{Y}
      \longrightarrow \Omega_Y \longrightarrow 0,
    \]
    a class $[e_{\Yt}] \in \Ext^{1}_{Y}(\Omega_Y, \cE)$.
  \item[\textup{(3)}] Two reflexive ribbons $\Yt$ and $\Yt'$ over
    $Y$ with conormal sheaf $\cE$ are isomorphic over $Y$ if and only
    if there exists $\alpha \in \Aut(\cE)$ such that
    $[e_{\Yt'}] = \alpha \cdot [e_{\Yt}]$.
\end{enumerate}
In particular, the set of isomorphism classes of reflexive ribbons
over $Y$ with conormal sheaf $\cE$ is in bijection with
$\Ext^{1}_{Y}(\Omega_Y, \cE) / \Aut(\cE)$.
\end{proposition}
 
\begin{remark}
When $\cE$ is a line bundle on a proper connected reduced scheme,
$\Aut(\cE) = k^{*}$, and the quotient
$\Ext^1(\Omega_Y, \cE)/k^*$ is a projective space (or empty); this
recovers the classical result of Bayer--Eisenbud
\cite[Corollary~1.4]{BE95}. For a reflexive sheaf of higher rank,
$\Aut(\cE)$ can be a more complex algebraic group, and the orbit
space need not have such a simple structure.
\end{remark}

\bibliographystyle{plain}
\begin{thebibliography}{KMM92}
\color{black}



\bibitem{AHL10}
Artebani, Michela; Hausen, Jürgen; Laface, Antonio
\emph{On Cox rings of K3 surfaces}
Compos. Math. 146 (2010), no. 4, 964–998.

\bibitem{ABS17}
Arbarello, Enrico; Bruno, Andrea; Sernesi, Edoardo
\emph{On hyperplane sections of K3 surfaces.}
Algebr. Geom. 4 (2017), no. 5, 562–596.


\bibitem{AltmannStevens}
K.~Altmann and J.~Stevens,
\newblock \emph{Cotangent cohomology of rational surface singularities},
\newblock Invent. Math. \textbf{138} (1999), no.~1, 163--181.

\bibitem{BE95}
Bayer, Dave; Eisenbud, David
\emph{Ribbons and their canonical embeddings}
Trans. Amer. Math. Soc. 347 (1995), no. 3, 719–756.


\bibitem{Bea02}
Beauville, A. \emph{Fano threefolds and $K3$ surfaces}
https://arxiv.org/abs/math/0211313
%\bibitem[Bel24]{Bel24}
%Belmans, Peter
%\emph{Fanography: a tool to visually study the geography of Fano 3-folds.}
%https://www.fanography.info

\bibitem{BGM24} Bangere, Purnaprajna; Gallego, Francisco Javier; Mukherjee; Jayan \emph{Projective smoothing of varieties with simple normal crossings.}(https://arxiv.org/abs/2403.04167)

\bibitem{BKM25}
Bayer, Arend;  Kuznetsov, Alexander;  Macr\'i, Emanuele 
\emph{Mukai models of Fano varieties}
arXiv:2501.16157

\bibitem{BM87}
Beauville, A.; Mérindol, J.-Y.
\emph{Sections hyperplanes des surfaces $K3$.[Hyperplane sections of $K3$ surfaces]}
Duke Math. J. 55 (1987), no. 4, 873–878.

\bibitem{BMR21}
Bangere, Purnaprajna; Mukherjee, Jayan; Raychaudhury, Debaditya
\emph{K3 carpets on minimal rational surfaces and their smoothings.}
Internat. J. Math.32(2021), no.6, Paper No. 2150032, 20 pp.

\bibitem{BM24} 
Bangere, Purnaprajna; Mukherjee, Jayan  
\emph{Extendability of projective varieties via degeneration to ribbons with applications to Calabi-Yau threefolds} (arXiv:2409.03960)

\bibitem{BM25}
Bangere, Purnaprajna; Mukherjee, Jayan  
\emph{Extendability of general prime $K3$ surfaces via degeneration to ribbons} (in progress)



%\bibitem[BGM24]{BGM24} Bangere, Purnaprajna; Gallego, Francisco Javier; Mukherjee; Jayan \emph{Projective degenerations of $K3$ surfaces to a union of Hirzebruch surfaces not embedded as scrolls.} (in preparation)

%\bibitem[BGM24b]{BGM24b} Bangere, Purnaprajna; Gallego, Francisco Javier; Mukherjee; Jayan \emph{Projective smoothing of simple normal crossing projective bundles.} (in preparation)

%\bibitem[BGM24c]{BGM24c} Bangere, Purnaprajna; Gallego, Francisco Javier; Mukherjee; Jayan \emph{Projective smoothing of simple normal crossing complete intersections.} (in preparation)

\bibitem{CD22}
Ciliberto, Ciro; Dedieu, Thomas
\emph{$K3$ curves with index $k>1$.}
Boll. Unione Mat. Ital. 15 (2022), no. 1-2, 87–115.

\bibitem{CD24}
Ciliberto, Ciro; Dedieu, Thomas
\emph{Double covers and extensions.}
Kyoto J. Math. 64 (2024), no. 1, 75–94.


\bibitem{CDGK20}
Ciliberto, C; Dedieu, T;  Galati, C.; Knutsen, A. L.  
\emph{Moduli of curves on Enriques surfaces.}
Adv. Math. 365 (2020), 107010, 42 pp. 6, 13, 14

\bibitem{CDS20}
Ciliberto, Ciro; Dedieu, Thomas; Sernesi, Edoardo
\emph{Wahl maps and extensions of canonical curves and $K3$ surfaces.}
J. Reine Angew. Math. 761 (2020), 219–245.

\bibitem{CLM93}
Ciliberto, C; Lopez, A; Miranda, R
\emph{Projective degenerations of $K3$ surfaces, Gaussian maps and Fano threefolds}
Invent. Math. 114, 641 667 (1993)

\bibitem{CLM94}
Ciliberto, Ciro; Lopez, Angelo Felice; Miranda, Rick
\emph{On the corank of Gaussian maps for general embedded $K3$ surfaces}
arXiv:alg-geom/9411008

\bibitem{CLM98}
Ciliberto, Ciro; Lopez, Angelo Felice; Miranda, Rick
\emph{Classification of varieties with canonical curve section via Gaussian maps on canonical curves.}
Amer. J. Math.   120 (1998), no. 1, 1–21.

\bibitem{CLM98c}
Ciliberto, Ciro; Lopez, Angelo Felice; Miranda, Rick
\emph{Corrigendum to ``Classification of varieties with canonical curve section via Gaussian maps on canonical curves''.}
Amer. J. Math. 143 (2021), no. 6, 1661–1663.

\bibitem{CilibertoLopezMiranda}
C.~Ciliberto, A.~F.~Lopez, and R.~Miranda,
\newblock \emph{Some remarks on the obstructedness of cones over curves of low genus},
\newblock in \emph{Higher Dimensional Complex Varieties (Trento, 1994)},
de Gruyter, Berlin, 1996, pp.~167--182.

\bibitem{CM90}
Ciliberto, Ciro; Miranda, Rick
\emph{On the Gaussian map for canonical curves of low genus.}
Duke Math. J. 61 (1990), no. 2, 417–443.

\bibitem{DM92}
Duflot, Jeanne; Miranda, Rick
\emph{The Gaussian map for rational ruled surfaces.}
Trans. Amer. Math. Soc. 330 (1992), no. 1, 447–459.

\bibitem{DM24}
Deopurkar, Anand; Mukherjee, Jayan
\emph{Syzygies of canonical ribbons on higher genus curves}
arXiv:2412.05500

%\color{black}
%\bibitem[Deb01]{Debarre} Debarre, O., \emph{Higher-dimensional algebraic geometry}, 
%Universitext Springer--Verlag, New York, 2001.
%\color{black}

\bibitem{DP12}
Duflot, Jeanne; Peters, Pamela L.
\emph{Gaussian maps for double covers of toric surfaces.}
Rocky Mountain J. Math. 42 (2012), no. 5, 1471–1520.












%\bibitem[Fri83]{F83}
%Friedman, Robert; 
%\emph{Global smoothings of varieties with normal crossings.} 
%Ann. of Math. (2) 118 (1983), no. 1, 75--114.

 

\color{black}

%\bibitem[Fuj14a]{KFujita} Fujita, Kento, \emph{Simple normal crossing Fano varieties and log Fano manifolds.}
%Nagoya Math. J. 214 (2014), 95-123.

%\bibitem[Fuj14b]{Fujino} Fujino, Osamu.
%\emph{Fundamental theorems for semi log canonical pairs.} Algebr. Geom. 1 (2014), no.2, 194-228.

%\color{black}

\bibitem{GGP08}
Gallego, Francisco Javier; González, Miguel; Purnaprajna, Bangere P., 
\emph{$K3$ double structures on Enriques surfaces and their smoothings.}
J. Pure Appl. Algebra 212 (2008), no. 5, 981--993

\bibitem{GGP13}
Gallego, F.J., González, M., Purnaprajna, B.P., 
\emph{An infinitesimal condition to smooth ropes.}
Rev Mat Complut 26, 253–269 (2013).

\bibitem{Gon06}
  Gonz\'alez, Miguel.
  \emph{Smoothing of ribbons over curves}.
  J. Reine Angew. Math. 591 (2006), 201--235.

\bibitem{GP97}
Gallego, F.J., Purnaprajna, B.P.
\emph{Degenerations of $K3$ surfaces in projective space}
Transactions of the American Mathematical Society
Volume 349, Number 6, June 1997, Pages 2477–2492



\color{black}

  %\bibitem[Gro61]{EGA} A. Grothendieck, \emph{\'El\'ements 
%de g\'eom\'etrie 
%alg\'ebrique. III. \'Etude cohomologique des faisceaux 
%coh\'erents. I.}, Inst. 
%Hautes \'Etudes Sci. Publ. Math.  
%{11}, (1961).

\bibitem{Hartshorne}  
Hartshorne, R; \emph{Algebraic geometry}, 
Grad. Texts in Math., No. 52
Springer-Verlag, New York-Heidelberg, 1977

\color{black}

\bibitem{Hor74}
Horikawa, Eiji
\emph{On deformations of holomorphic maps. II.}
J. Math. Soc. Japan 26 (1974), 647–667.



\bibitem{I77}
Iskovskih, V. A.
\emph{Fano threefolds. I.}
(Russian)Izv. Akad. Nauk SSSR Ser. Mat.41 (1977), no.3, 516–562, 717.

\bibitem{I78}
Iskovskih, V. A.
\emph{Fano threefolds. II.}
(Russian)Izv. Akad. Nauk SSSR Ser. Mat.42 (1978), no.3, 506–549.



%\bibitem[Kim17]{Ki17}
%Kimura, Y.;
%\emph{Structure of stable degeneration of $K3$ surfaces into pairs of rational elliptic surfaces}
%JHEP03(2018)045, %https://doi.org/10.48550/arXiv.1710.04984
\bibitem{K01}
Knutsen, Andreas Leopold
\emph{On $k$th-order embeddings of $K3$ surfaces and Enriques surfaces.}
Manuscripta Math. 104 (2001), no. 2, 211–237.

\bibitem{K20}
Knutsen, Andreas Leopold
\emph{Global sections of twisted normal bundles of $K3$ surfaces and their hyperplane sections.}
Atti Accad. Naz. Lincei Rend. Lincei Mat. Appl. 31 (2020), no. 1, 57–79.


\bibitem{KLM11}
Knutsen, Andreas Leopold; Lopez, Angelo Felice; Muñoz, Roberto
\emph{On the extendability of projective surfaces and a genus bound for Enriques-Fano threefolds.}
J. Differential Geom.88(2011), no.3, 485–518.

\bibitem{KL81}
Kleppe, J.; 
\emph{The Hilbert Flag Scheme, its properties, and its connection with the Hilbert Scheme. Applications to curves in three-space. Thesis. University of Oslo, Preprint No. 5 (1981)}
\bibitem{KL87}
Kleppe, J.;
\emph{Non-reduced components of the Hilbert scheme of smooth space curves. In: Space Curve.}
Rocca di papa, 1985. Ghione, F., Peskine, C., Sernesi, E. (eds.) Springer Lect. Notes Math. No. 1266 (1987), 181-207

\bibitem{HR} M.~Hochster and J.\,L.~Roberts, \emph{Rings of invariants of reductive groups acting on regular rings are Cohen--Macaulay}, Advances in Math.\ \textbf{13} (1974), 115--175.
 
\bibitem{Illusie} L.~Illusie, \emph{Complexe cotangent et d\'eformations I, II}, Lecture Notes in Math.\ \textbf{239}, \textbf{283}, Springer-Verlag, 1971--1972.
 
\bibitem{Saito} M.~Saito, \emph{Mixed Hodge complexes on algebraic varieties}, Math.\ Ann.\ \textbf{316} (2000), 283--331.

\bibitem{stacks-project}
The Stacks Project Authors,
\emph{The Stacks Project},
\url{https://stacks.math.columbia.edu}
 
\bibitem{Steenbrink} J.\,H.\,M.~Steenbrink, \emph{Mixed Hodge structure on the vanishing cohomology}, in: Real and Complex Singularities, Sijthoff and Noordhoff, 1977, 525--563.
 
\bibitem{Wehler} J.~Wehler, \emph{Cyclic coverings: deformation and Torelli theorem}, Math.\ Ann.\ \textbf{274} (1986), 443--472.

  




%\bibitem[Kon85]{Ko85}
%Kondo, S
%\emph{Type II degenerations of $K3$ surfaces}
%Nagoya Math. J. Vol. 99 (1985), 11-30



%\bibitem[Kul77]{Ku77}
%Kulikov, V. 
%\emph{Degenerations of K3 surfaces and Enriques surfaces} Math. USSR Izvestiya 11 (1977), 957-989.





\color{black}



\bibitem{Lop23}
Lopez, Angelo Felice
\emph{On the extendability of projective varieties: a survey.The art of doing algebraic geometry, 261–291. With an appendix by Thomas Dedieu}
Trends Math. Birkhäuser/Springer, Cham, [2023], ©2023 ISBN:978-3-031-11937-8
ISBN:978-3-031-11938-5


%\bibitem[LR12]{LR}
%Liu, Wenfei; Rollenske, Sönke. \emph{Two-dimensional semi-log-canonical hypersurfaces}, 
%Matematiche (Catania) 67 (2012), no. 2, 185–202.

\bibitem{Lvo92} 
L\'vovsky, S.
\emph{Extensions of projective varieties and deformations. I, II.}
Michigan Math. J.39(1992), no.1, 41–51, 65--70.





\bibitem{MM82}
Mori, Shigefumi; Mukai, Shigeru
\emph{Classification of Fano 3-folds with $B2 \geq 2$.}
Manuscripta Math.36(1981/82/1982), no.2, 147–162.


\bibitem{M88}
Mukai, S
\emph{Curves, K3 Surfaces and Fano 3-folds of Genus $\leq$ 10}
Algebraic Geometry and Commutative Algebra, Academic Press,
1988, Pages 357-377.

  \color{black}
  


  \color{black}
%\bibitem[Par91]{Par91}
%Pardini, Rita
%\emph{Abelian covers of algebraic varieties.}
%J. Reine Angew. Math.417(1991), 191–213.




%\bibitem[PP81]{PP81}
%Persson, U; Pinkham, H 
%\emph{Degenerations of surfaces with trivial canonical bundle}
%Ann. of Math. 113 (1981), 45-66

\bibitem{Pinkham}
H.~C.~Pinkham,
\newblock \emph{Deformations of algebraic varieties with \(G_m\) action},
\newblock Ast\'erisque \textbf{20}, Soci\'et\'e Math\'ematique de France, Paris, 1974.


\bibitem{Ser}
Sernesi, Edoardo.
\emph{Deformations of Algebraic Schemes.}
Grundlehren der Mathematischen Wissenschaften [Fundamental Principles of Mathematical Sciences], 334. Springer-Verlag, Berlin, 2006.

\bibitem[Stacks]{stacks-project}
The Stacks Project Authors,
\emph{Stacks Project}.
\url{https://stacks.math.columbia.edu}.

\bibitem{StacksProject}
The Stacks Project Authors,
\emph{Stacks Project},
\url{https://stacks.math.columbia.edu}.

\bibitem{Tot20}
Totaro, Burt
\emph{Bott vanishing for algebraic surfaces.}
Trans. Amer. Math. Soc. 373 (2020), no. 5, 3609–3626.
%\bibitem[Sei92]{Sei92}
%Seiler, Wolfgang K.
%\emph{Deformations of ruled surfaces.}
%J. Reine Angew. Math.426(1992), 203–219.

%\bibitem[SP]{SP}
%Stacks project authors
%\emph{The Stacks project}
%\url{https://stacks.math.columbia.edu}, 2023

%\bibitem[Tho13]{T13}
%Thompson, A;
%\emph{Degenerations of $K3$ surfaces of degree two}
%Transactions of the American Mathematical Society
%Volume 366, Number 1, January 2014, Pages 219–243
%S 0002-9947(2013)05759-5


\bibitem{V92}
Voisin, Claire
\emph{Sur l'application de Wahl des courbes satisfaisant la condition de Brill-Noether-Petri.[On the Wahl mapping of curves satisfying the Brill-Noether-Petri condition]}
Acta Math. 168 (1992), no. 3-4, 249–272.






\bibitem{W87}
Wahl, J. M. 
\emph{The Jacobian algebra of a graded Gorenstein singularity.}
Duke Math. J. 55 (1987), no. 4,
843-871

\bibitem{Zak91}
Zak, F. L.
\emph{Some properties of dual varieties and their applications in projective geometry.}
Algebraic geometry (Chicago, IL, 1989), 273–280.
Lecture Notes in Math., 1479
Springer-Verlag, Berlin, 1991

\bibitem{Wehler1986}
J.~Wehler,
\newblock Cyclic coverings: deformation and Torelli theorem,
\newblock \emph{Math. Ann.} \textbf{274} (1986), 443--472.

\bibitem{Gonzalez2005}
M.~Gonz\'alez,
\newblock Smoothing of ribbons over curves,
\newblock \emph{J. Reine Angew. Math.} \textbf{591} (2006), 201--235.

\bibitem{GallegoGonzalezPurnaprajna2010}
F.~J. Gallego, M.~Gonz\'alez, and B.~P. Purnaprajna,
\newblock Deformation of canonical morphisms and the moduli of surfaces of general type,
\newblock \emph{Invent. Math.} \textbf{182} (2010), 1--46.

\bibitem{Illusie1971}
L.~Illusie,
\newblock \emph{Complexe cotangent et d\'eformations I},
\newblock Lecture Notes in Mathematics, Vol.~239,
\newblock Springer-Verlag, Berlin--New York, 1971.

\bibitem{Sernesi2006}
E.~Sernesi,
\newblock \emph{Deformations of Algebraic Schemes},
\newblock Grundlehren der mathematischen Wissenschaften, Vol.~334,
\newblock Springer-Verlag, Berlin, 2006.




 
\end{thebibliography}

\end{document}






\begin{remark}[2.4]
The content of Lemma~2.3 is a purely cohomological fact of sheaves on $X$.
The proof given in \cite[4.2]{Hor74}, in the context of complex manifolds, is still valid
in the case of a homomorphism of quasi--coherent sheaves
\[
\mathcal A \xrightarrow{\,F\,} \mathcal B
\]
with kernel and cokernel, respectively, $\mathcal K$ and $\mathcal N$, on a noetherian separated scheme $X$.
Let $\mathcal V$ be an open cover of $X$. We set
\begin{equation}\tag{2.4.1}
D(X,F;\mathcal V)
=
\frac{\Bigl\{(g,\rho)\in C^{0}(\mathcal V,\mathcal B)\times Z^{1}(\mathcal V,\mathcal A)\ \Big|\ 
\delta g = F\rho\Bigr\}}
{\Bigl\{(Fh,\delta h)\ \Big|\ h\in C^{0}(\mathcal V,\mathcal A)\Bigr\}}
\end{equation}
and
\begin{equation}\tag{2.4.2}
D(X,F)=\varinjlim_{\mathcal V}\, D(X,F;\mathcal V),
\end{equation}
where the direct limit is taken under refinement of coverings. Then for every open affine cover
$\mathcal V$ the natural map $D(X,F;\mathcal V)\to D(X,F)$ is an isomorphism.
Two exact sequences, like in Lemma~2.3, are then obtained:
\[
H^{0}(\mathcal A)\longrightarrow H^{0}(\mathcal B)\longrightarrow D(X,F)
\longrightarrow H^{1}(\mathcal A)\longrightarrow H^{1}(\mathcal B),
\]
\[
0\longrightarrow H^{1}(\mathcal K)\longrightarrow D(X,F)\longrightarrow H^{0}(\mathcal N)
\longrightarrow H^{2}(\mathcal K).
\]
\end{remark}